% 先修：Lebesgue 积分、单调收敛、Fatou 引理、Tonelli/Fubini 定理。
\section{从 Young 不等式到 \texorpdfstring{$L^p$}{Lp} 空间的完备性}\label{sec:lp-foundations}

本专题服务于后面的流体弱解理论。学习主线是：先证明积分估计，再证明函数集合带范数，接着证明 Cauchy 序列在空间内有极限；最后把同样的方法用于函数及其弱导数。全文的 Sobolev 指索伯列夫空间，Young 指杨氏不等式，H\"older 指赫尔德不等式。

\subsection{预备知识、几乎处处与学习路线}

本专题使用 Lebesgue 积分。先承认以下测度积分定理：非负可测函数的单调收敛定理、Fatou 引理、Tonelli 定理及可积函数的 Fubini 定理。它们的测度论证明不在本专题内。记 $\Omega\subset\R^d$ 为非空开集，$|E|$ 为 Lebesgue 测度，$\|\cdot\|_p$ 默认在 $\Omega$ 上计算，$1\le p\le\infty$。

“几乎处处”（a.e.）表示可以去掉一个零测集。例如仅在一点改变函数，不会改变任何 $L^p$ 积分。以下以实值标量函数为主；向量场按欧氏长度定义 $L^p$ 范数，各分量弱求导。复值情形将内积中的第二因子取共轭即可。

\begin{keypoint}
逻辑依赖是 Young $\Rightarrow$ H\"older $\Rightarrow$ Minkowski $\Rightarrow L^p$ 为赋范空间 $\Rightarrow L^p$ 完备 $\Rightarrow W^{k,p}$ 完备。嵌入定理再把“导数的可积性”转化为“函数更好的可积性或连续性”。
\end{keypoint}

\subsection{Young 不等式：先证明一个数的乘积估计}

\begin{theorem}[Young 不等式]\label{thm:lp-young}
设 $1<p<\infty$，共轭指数 $p'=p/(p-1)$，即 $1/p+1/p'=1$。对 $a,b\ge0$，
\begin{equation}
 ab\le\frac{a^p}{p}+\frac{b^{p'}}{p'},
 \qquad\text{等号当且仅当 }a^p=b^{p'}.
 \label{eq:lp-young}
\end{equation}
\end{theorem}
\begin{proof}
固定 $b$，考虑 $F(a)=a^p/p-ab+b^{p'}/p'$。其导数为 $F'(a)=a^{p-1}-b$，在 $a=b^{1/(p-1)}$ 处取得全局最小值。代入并用 $1/p+1/p'=1$ 得最小值为零；$b=0$ 的情形直接成立。这也给出等号条件。
\end{proof}

将 $a,b$ 换为 $\lambda a,b/\lambda$，可得常用于吸收项的版本：对任意 $\varepsilon>0$，
\begin{equation}
 ab\le\varepsilon a^p+C_{p,\varepsilon}b^{p'},\qquad
 C_{p,\varepsilon}=\frac{(p\varepsilon)^{-p'/p}}{p'}.
 \label{eq:lp-young-eps}
\end{equation}
尤其 $ab\le\varepsilon a^2+b^2/(4\varepsilon)$。能量估计中，选 $\varepsilon$ 足够小，就能把第一项并入左侧的黏性耗散。

\subsection{定义 \texorpdfstring{$L^p$}{Lp}：为什么必须取等价类}

先定义 $\mathcal L^p(\Omega)=\{f\text{ 可测}:\int_\Omega|f|^p\dd x<\infty\}$（$p<\infty$）。规定 $f\sim g$ 当且仅当 $f=g$ a.e.，取商空间
\begin{equation}
 L^p(\Omega)=\mathcal L^p(\Omega)/\!\sim,\qquad
 \|f\|_p=\left(\int_\Omega|f|^p\dd x\right)^{1/p}.
 \label{eq:lp-definition}
\end{equation}
定义时还不能直接宣称它是范数，三角不等式将在后面证明。若不取等价类，函数 $\mathbf1_{\{x_0\}}$ 逐点并非零函数，却有积分“长度”零，只能得到半范数。

在 $p=\infty$ 时定义
\begin{equation}
 \|f\|_\infty=\operatorname*{ess\,sup}_{x\in\Omega}|f(x)|
 =\inf\{M\ge0:|f|\le M\text{ a.e.}\},
 \label{eq:lp-infinity}
\end{equation}
$L^\infty$ 是本质有界函数的 a.e. 等价类。例如 $f(1/n)=n$ 而其余点取零，普通上确界为无穷，本质上确界为零。$p=\infty$ 的符号不能解释成直接把无穷代入积分公式。

\subsection{H\"older 不等式：从逐点估计到积分估计}

\begin{theorem}[H\"older 不等式]\label{thm:lp-holder}
若 $f\in L^p(\Omega)$、$g\in L^{p'}(\Omega)$，则 $fg\in L^1(\Omega)$ 且
\begin{equation}
 \boxed{\int_\Omega|fg|\dd x\le\|f\|_p\|g\|_{p'}.}
 \label{eq:lp-holder}
\end{equation}
这里允许端点 $(p,p')=(1,\infty)$ 或 $(\infty,1)$。
\end{theorem}
\begin{proof}
先设 $1<p<\infty$。任一范数为零时结论直接成立；否则令 $F=|f|/\|f\|_p$、$G=|g|/\|g\|_{p'}$。逐点应用 Young，再积分：
\[
 \int_\Omega FG\dd x
 \le\frac1p\int_\Omega F^p\dd x+
 \frac1{p'}\int_\Omega G^{p'}\dd x=1.
\]
乘回两个范数即得结论。在端点处用 $|fg|\le|f|\|g\|_\infty$ a.e. 后积分。
\end{proof}

非零函数且 $1<p<\infty$ 时，等号条件是 $|f|^p/\|f\|_p^p=|g|^{p'}/\|g\|_{p'}^{p'}$ a.e.。特别 $p=2$ 给出 Cauchy--Schwarz 不等式。

反复使用 H\"older 可得多因子版本：若 $\sum_j1/p_j=1$，则 $\int\prod_j|f_j|\le\prod_j\|f_j\|_{p_j}$；若 $1/r=1/p+1/q$ 且 $r\ge1$，则
\begin{equation}
 \|fg\|_r\le\|f\|_p\|g\|_q.
 \label{eq:lp-product}
\end{equation}
证明后式时，对 $|f|^r|g|^r$ 用指数 $p/r,q/r$ 的 H\"older，再取 $r$ 次根。后面三维对流项使用的指数是 $6,2,3$，因为 $1/6+1/2+1/3=1$。

\subsection{Minkowski 不等式与范数三公理}

\begin{theorem}[Minkowski 不等式]\label{thm:lp-minkowski}
对 $f,g\in L^p(\Omega)$，有 $f+g\in L^p(\Omega)$，并且
\begin{equation}
 \boxed{\|f+g\|_p\le\|f\|_p+\|g\|_p.}
 \label{eq:lp-minkowski}
\end{equation}
\end{theorem}
\begin{proof}
当 $1<p<\infty$，先由数函数 $s\mapsto s^p$ 的凸性得到
$|f+g|^p\le2^{p-1}(|f|^p+|g|^p)$，所以 $f+g$ 已确定属于 $L^p$，避免后面除以一个未经证明有限的量。由 H\"older，
\begin{align*}
 \|f+g\|_p^p
 &\le\int_\Omega (|f|+|g|)|f+g|^{p-1}\dd x\\
 &\le(\|f\|_p+\|g\|_p)
 \big\||f+g|^{p-1}\big\|_{p'}\\
 &=(\|f\|_p+\|g\|_p)\|f+g\|_p^{p-1}.
\end{align*}
若 $\|f+g\|_p=0$ 则已成立，否则相除。$p=1$ 直接积分逐点三角不等式；$p=\infty$ 对 a.e. 不等式取本质上确界。
\end{proof}

现在逐条核对：
\begin{enumerate}
 \item 非负及正定性：$\|f\|_p\ge0$，且 $\|f\|_p=0\iff f=0$ a.e.，即商空间中的零元素；
 \item 绝对齐次性：$\|cf\|_p=|c|\|f\|_p$；
 \item 三角不等式：正是 Minkowski。
\end{enumerate}
加法、数乘对 a.e. 等价类良定义，由上面的可积性和齐次性封闭。因此 $L^p(\Omega)$ 是赋范线性空间。

\begin{example}[为什么通常限制 $p\ge1$]
在 $(0,2)$ 取 $f=\mathbf1_{(0,1)}$、$g=\mathbf1_{(1,2)}$。当 $0<p<1$，$\|f+g\|_p=2^{1/p}>2=\|f\|_p+\|g\|_p$，三角不等式失效。此时仍可研究 $L^p$，但这里的表达式不是范数。
\end{example}

\subsection{完备性：每个 Cauchy 序列都在空间内收敛}

\begin{definition}[Cauchy、完备与 Banach 空间]
若任意 $\varepsilon>0$ 都存在 $N$，使 $m,n\ge N$ 时 $\|f_n-f_m\|_p<\varepsilon$，称序列为 Cauchy。若每个这样的序列都有 $f\in L^p$ 满足 $\|f_n-f\|_p\to0$，空间称为完备；完备赋范空间称为 Banach 空间。
\end{definition}

\begin{theorem}[$L^p$ 完备性]\label{thm:lp-complete}
对任意 $1\le p\le\infty$，$L^p(\Omega)$ 为 Banach 空间。$p=2$ 时还是内积 $\langle f,g\rangle=\int fg$ 诱导的 Hilbert 空间。
\end{theorem}
\begin{proof}
先设 $p<\infty$，令 $(f_n)$ 为 Cauchy 序列。

\textbf{第一步：抽取相邻差可求和的子列。}选递增的 $n_k$，使
\[
 \|f_{n_{k+1}}-f_{n_k}\|_p\le2^{-k}.
\]
令 $h_k=|f_{n_{k+1}}-f_{n_k}|$，$S_m=\sum_{k=1}^mh_k$。有限次 Minkowski 给出 $\|S_m\|_p\le\sum_{k=1}^m2^{-k}\le1$。

\textbf{第二步：构造逐点极限，并证明它可积。}由于 $S_m\uparrow S=\sum_{k=1}^\infty h_k$，单调收敛定理给出 $\int S^p\le1$。所以 $S<\infty$ a.e.，且 $S\in L^p$。在这个满测集上，数列的差绝对可求和，故可定义
\[
 f=f_{n_1}+\sum_{k=1}^\infty(f_{n_{k+1}}-f_{n_k}).
\]
在例外零测集上任取零。$f$ 可测，且 $|f|\le|f_{n_1}|+S$，因此 $f\in L^p$。

\textbf{第三步：从逐点收敛加强为范数收敛。}对 $T_k=\sum_{j=k}^\infty h_j$，再次用有限和的 Minkowski 与单调收敛，
\[
 \|f-f_{n_k}\|_p\le\|T_k\|_p
 \le\sum_{j=k}^\infty2^{-j}=2^{1-k}\longrightarrow0.
\]

\textbf{第四步：恢复整个序列。}给定 $\varepsilon>0$，先令 Cauchy 尾部距离小于 $\varepsilon/2$，再取尾部子列元素使 $\|f_{n_k}-f\|_p<\varepsilon/2$。三角不等式给出 $\|f_n-f\|_p<\varepsilon$。

最后处理 $p=\infty$。可取同样的子列。对每个 $k$，相邻差的逐点界 $|f_{n_{k+1}}-f_{n_k}|\le2^{-k}$ 在某个满测集上成立；再纳入 $f_{n_1}$ 的本质有界例外集。去掉这些\textbf{可数个}零测集的并后，级数在剩余集合上一致收敛，得到本质有界 $f$，且 $\|f-f_{n_k}\|_\infty\le2^{1-k}$。同样恢复全序列。$p=2$ 的内积性质直接由积分与 Cauchy--Schwarz 得到，完备性已经证明。
\end{proof}

\begin{warningbox}
Cauchy 并不表示逐点单调，也不能未经论证直接写出逐点极限。证明的关键是\textbf{先抽取差范数可求和的子列}。完备性也不是说所有有界序列都收敛。
\end{warningbox}

\subsection{包含关系、插值与收敛的区别}

若 $|\Omega|<\infty$ 且 $1\le p<q\le\infty$，对 $|f|^p\cdot1$ 使用 H\"older 得
\begin{equation}
 \|f\|_p\le|\Omega|^{1/p-1/q}\|f\|_q,
 \qquad L^q(\Omega)\hookrightarrow L^p(\Omega).
 \label{eq:lp-finite-inclusion}
\end{equation}
在无限测度区域上不能照搬，例如 $1\notin L^p(\R^d)$（$p<\infty$），但 $1\in L^\infty$。$(0,1)$ 上 $x^{-\alpha}\in L^p$ 当且仅当 $\alpha p<1$，说明有限测度下的包含一般也不可反向。

若 $1\le a<b\le\infty$，$0\le\theta\le1$ 且 $1/r=\theta/a+(1-\theta)/b$，则
\begin{equation}
 \|f\|_r\le\|f\|_a^\theta\|f\|_b^{1-\theta}.
 \label{eq:lp-interpolation}
\end{equation}
当指数有限时，把 $|f|^r$ 分成 $|f|^{\theta r}|f|^{(1-\theta)r}$，使用 H\"older 的指数 $a/(\theta r)$、$b/((1-\theta)r)$；$b=\infty$ 时直接提出本质上界。这是后面估计 $L^3,L^4$ 范数的工具。

对于 $p<\infty$，范数收敛 $f_n\to f$ 推出依测度收敛，因为
\[
 |\{|f_n-f|>\varepsilon\}|\le\varepsilon^{-p}\|f_n-f\|_p^p.
\]
也能抽出 a.e. 收敛子列；可选 $\|f_{n_k}-f\|_p^p\le2^{-k}$，利用 Tonelli 得 $\sum_k|f_{n_k}-f|^p<\infty$ a.e.。但全序列不必逐点收敛。反过来，$f_n=n^{1/p}\mathbf1_{(0,1/n)}$ a.e. 趋零而 $\|f_n\|_p=1$，说明逐点收敛本身不足以推出范数收敛。

当 $1<p<\infty$ 时，弱收敛写为 $f_n\rightharpoonup f$，可用 $\int f_ng\to\int fg$ 对所有 $g\in L^{p'}$ 检验。$L^p$ 的对偶表示与自反性属于进一步的泛函分析定理，这里陈述使用而不证明。自反性保证有界序列存在弱收敛子列；强收敛子列还需要紧性。

% 由 lp-sobolev.tex 载入；先修为 Lebesgue 积分的基础定理。
\clearpage
\section{弱导数与 Sobolev 空间：从分部积分逐步建立理论}\label{sec:sobolev-definition}

本节不预先把“弱导数”当作会计算的符号，而是依次回答：为什么这样定义，能否唯一确定，何时存在，如何用光滑函数逼近，以及边界条件如何进入空间。凡出现 $D_i u$，都应能还原为一个对所有测试函数成立的积分恒等式。“弱导数”描述求导的意义；“弱收敛”描述序列的收敛方式，是两个不同概念。

\subsection{从经典分部积分推出弱导数的定义}

先设 $u\in C^1(\Omega)$，$\varphi\in C_c^\infty(\Omega)$。由于 $\varphi$ 支撑远离边界，$u\varphi$ 也紧支撑，故
\[
 0=\int_\Omega\partial_i(u\varphi)\dd x
  =\int_\Omega(\partial_i u)\varphi\dd x
   +\int_\Omega u\partial_i\varphi\dd x.
\]
移项得 $\int u\partial_i\varphi=-\int(\partial_i u)\varphi$。右式中导数已从未知函数移到已知的光滑测试函数，因此即使 $u$ 不经典可微，左侧仍然有定义。

\begin{definition}[测试函数、局部可积与弱导数]\label{def:weak-derivative}
$u\in L^1_{\mathrm{loc}}(\Omega)$ 意味着在每个紧子集上可积。若存在 $g_i\in L^1_{\mathrm{loc}}(\Omega)$，使
\begin{equation}
 \boxed{\int_\Omega u\partial_i\varphi\dd x
       =-\int_\Omega g_i\varphi\dd x,
       \quad\forall\varphi\in C_c^\infty(\Omega),}
 \label{eq:weak-derivative-definition}
\end{equation}
则称 $g_i$ 为 $u$ 的第 $i$ 个弱偏导数，记为 $D_i u$。
\end{definition}

局部可积性保证所有积分有限：$\varphi$ 及其导数有界且支撑在同一个紧集内。定义中的“任意 $\varphi$”至关重要；只验证一两个测试函数远远不够。测试函数紧支撑消除了边界项，\textbf{并未要求 $u$ 的边界值为零}。

对多重指标 $\alpha=(\alpha_1,\ldots,\alpha_d)$，记 $|\alpha|=\sum_i\alpha_i$，定义
\begin{equation}
 \int_\Omega uD^\alpha\varphi\dd x
 =(-1)^{|\alpha|}\int_\Omega D^\alpha u\,\varphi\dd x.
 \label{eq:weak-higher-derivative}
\end{equation}
每移走一个导数就引入一个负号。重复求导与直接用此式定义一致；混合分布导数交换，因为测试函数的混合经典导数交换。

\subsection{先证明光滑核逼近：后续唯一性与运算法则的基础}

取非负光滑函数
\[
 \rho(x)=c\begin{cases}\exp\big(-1/(1-|x|^2)\big),&|x|<1,\\0,&|x|\ge1,\end{cases}
 \qquad\int_{\R^d}\rho=1,
\]
并令 $\rho_\varepsilon(x)=\varepsilon^{-d}\rho(x/\varepsilon)$。这是把一个总质量为 1 的光滑核压缩进半径 $\varepsilon$ 的球。

\begin{lemma}[平移连续与近似恒等]\label{lem:weak-approx-identity}
若 $f\in L^p(\R^d)$，$1\le p<\infty$，则
\[
 \|f(\cdot-h)-f\|_p\to0\quad(h\to0),\qquad
 \|\rho_\varepsilon*f-f\|_p\to0\quad(\varepsilon\to0).
\]
若 $f\in L^1_{\mathrm{loc}}(\Omega)$，第二结论在每个 $K\Subset\Omega$ 上以 $L^1(K)$ 成立。
\end{lemma}
\begin{proof}
\textbf{先证明平移。}给定 $\eta>0$，利用 $C_c(\R^d)$ 在 $L^p$ 中的稠密性，选 $g$ 使 $\|f-g\|_p<\eta$。这一基础稠密性可由简单函数逼近、Lebesgue 测度的内外正则性和连续截断证明，此处作为测度论预备使用。平移不改变 $L^p$ 范数，所以
\[
 \|f(\cdot-h)-f\|_p\le2\eta+\|g(\cdot-h)-g\|_p.
\]
$g$ 一致连续，小平移下两函数的支撑落在固定有限测度集内，故最后一项趋零。令 $\eta\downarrow0$ 得结论。

\textbf{再证明卷积。}由 Jensen 不等式和 Tonelli，
\begin{align*}
 \|\rho_\varepsilon*f-f\|_p^p
 &\le\int\rho_\varepsilon(h)
             \|f(\cdot-h)-f\|_p^p\dd h\\
 &\le\sup_{|h|\le\varepsilon}\|f(\cdot-h)-f\|_p^p\longrightarrow0.
\end{align*}
此处 Jensen 只是对凸函数 $t\mapsto|t|^p$ 和概率权重 $\rho_\varepsilon$ 使用凸性。局部情形取 $\chi\in C_c^\infty(\Omega)$，使它在 $K$ 的一个邻域上为 1，将 $\chi f$ 零延拓到全空间。充分小的 $\varepsilon$ 下，它与 $f$ 的卷积在 $K$ 上相同。
\end{proof}

对局部可积 $f$，$f_\varepsilon=\rho_\varepsilon*f$ 在内缩区域上光滑；对核逐次求导即可得到
\[
 D^\alpha f_\varepsilon(x)=\int f(y)D_x^\alpha\rho_\varepsilon(x-y)\dd y.
\]
在固定紧集上，核及其导数有共同紧支撑并有界，故可由控制收敛交换求导和积分。

\subsection{基本引理、弱导数唯一性和零梯度判据}

\begin{lemma}[局部可积函数的基本引理]\label{lem:weak-fundamental}
若 $h\in L^1_{\mathrm{loc}}(\Omega)$，且 $\int h\varphi=0$ 对所有 $\varphi\in C_c^\infty(\Omega)$ 成立，则 $h=0$ a.e.。
\end{lemma}
\begin{proof}
固定 $K\Subset\Omega$，并取 $\varepsilon<\operatorname{dist}(K,\partial\Omega)$。对每个 $x\in K$，函数 $y\mapsto\rho_\varepsilon(x-y)$ 是允许的测试函数，故 $(h*\rho_\varepsilon)(x)=0$。由引理~\ref{lem:weak-approx-identity}，它在 $L^1(K)$ 中趋于 $h$，所以 $\int_K|h|=0$。用可数个这样的紧集穷竭 $\Omega$，便得全区域 a.e. 为零。
\end{proof}

若 $g_i,\widetilde g_i$ 都满足弱导数定义，作差得 $\int(g_i-\widetilde g_i)\varphi=0$，基本引理立即给出 $g_i=\widetilde g_i$ a.e.。因此弱导数唯一到零测集；修改 $u$ 的零测集取值也不会改变它。

\begin{proposition}[连通区域上零弱梯度意味着常数]\label{prop:weak-zero-gradient}
若 $\Omega$ 开且连通，$u\in L^1_{\mathrm{loc}}(\Omega)$、$D_i u=0$ 对所有 $i$ 成立，则 $u$ a.e. 等于一个常数。
\end{proposition}
\begin{proof}
在球 $B\Subset\Omega$ 内作局部卷积。由定义，$\partial_i u_\varepsilon=0$，故在略小的同心球上 $u_\varepsilon=c_\varepsilon$。固定一个更小球 $B'$，局部 $L^1$ 收敛说明这些常数在 $L^1(B')$ 中 Cauchy，即
$|B'||c_\varepsilon-c_\delta|\to0$。所以 $c_\varepsilon\to c$，$u=c$ a.e. 于 $B'$。重叠球上的常数必须相同；开连通集内任意两点可由有限条相交小球连接，故整个区域的常数一致。
\end{proof}
若区域不连通，不同连通分支可以有不同常数。这正是零均值 Poincar\'e 不等式需要连通性的原因。

\subsection{分布导数总存在，函数型弱导数却未必存在}

分布是测试函数上的连续线性泛函。这里的连续性意为：若 $\varphi_n$ 支撑在同一个紧集内，且任意阶导数都一致趋零，则 $T(\varphi_n)\to0$。每个 $u\in L^1_{\mathrm{loc}}$ 产生分布 $T_u(\varphi)=\int u\varphi$，因为在固定紧集 $K$ 上
\[
 |T_u(\varphi)|\le\|u\|_{L^1(K)}\|\varphi\|_\infty.
\]
定义 $D_iT$ 的作用为 $\langle D_iT,\varphi\rangle=-\langle T,\partial_i\varphi\rangle$；连续性和线性仍然成立，所以\textbf{每个分布都可以任意次求分布导数}。但只有当 $D_iT_u=T_g$、$g\in L^1_{\mathrm{loc}}$ 时，才得到本专题的函数型弱导数。

例如 Dirac 分布 $\delta_{x_0}(\varphi)=\varphi(x_0)$ 是连续线性泛函，却不能由体积局部可积函数表示。弱导数的存在问题，实质上是检查分布导数是否具有所需可积性的函数表示。基本引理说明这种表示若存在则 a.e. 唯一。

\subsection{折角、跳跃、奇性与界面：逐项计算弱导数}

\begin{example}[折角：$u(x)=|x|$]
在 $(-1,1)$ 上，测试函数在端点附近为零。逐段积分得到
\begin{align*}
 \int_{-1}^1|x|\varphi'\dd x
 &=\int_{-1}^0(-x)\varphi'\dd x+\int_0^1x\varphi'\dd x\\
 &=[-x\varphi]_{-1}^0+\int_{-1}^0\varphi
   +[x\varphi]_0^1-\int_0^1\varphi\\
 &=-\int_{-1}^1\operatorname{sgn}(x)\varphi(x)\dd x.
\end{align*}
因此 $D|x|=\operatorname{sgn}x$ a.e.。再求一次分布导数：
$\langle D^2|x|,\varphi\rangle=-\int\operatorname{sgn}x\,\varphi'=2\varphi(0)$，所以 $D^2|x|=2\delta_0$。结论是 $|x|\in W^{1,\infty}$，但 $|x|\notin W^{2,1}$。
\end{example}

\begin{example}[跳跃：阶跃函数]
令 $H=\mathbf1_{(0,1)}$，定义域为 $(-1,1)$。由于 $\int H\varphi'=\int_0^1\varphi'=-\varphi(0)$，其分布导数为 $DH=\delta_0$。假设存在局部可积 $g$ 表示它；取 $0\le\varphi_\varepsilon\le1$、$\varphi_\varepsilon(0)=1$，支撑缩进 $(-\varepsilon,\varepsilon)$，则
\[
 1=\int g\varphi_\varepsilon
 \le\left|\int g\varphi_\varepsilon\right|
 \le\int_{-\varepsilon}^{\varepsilon}|g|\longrightarrow0,
\]
矛盾。最后一步来自可积函数积分的绝对连续性。因此 $H$ 不属于任何 $W^{1,p}(-1,1)$，即使它在原点以外的经典导数都为零。
\end{example}

\begin{example}[可积奇性：径向幂函数]
在 $B_1\subset\R^d$、$d\ge2$ 中取 $u(x)=|x|^{-a}$，$a>0$。极坐标计算给出
\[
 \int_{B_1}|u|^p=|\mathbb S^{d-1}|\int_0^1r^{d-1-ap}\dd r,
 \quad |\nabla u|=a r^{-a-1}.
\]
因此候选弱导数可积到 $p$ 次幂的条件是 $(a+1)p<d$。还须检查原点没有额外 Dirac 项：在 $B_1\setminus\overline B_\varepsilon$ 分部积分，内球边界项绝对值不超过 $C\varepsilon^{d-1-a}$，在上述条件下趋零。令 $\varepsilon\downarrow0$ 得到真正的弱导数 $D_i u=-a x_i|x|^{-a-2}$。反过来，若属于 $W^{1,p}$，在去掉原点的区域弱导数必须等于经典导数，故同一积分条件必要。于是
\[
 |x|^{-a}\in W^{1,p}(B_1)\iff (a+1)p<d\quad(a>0).
\]
三维中 $0<a<1/2$ 给出无界的 $H^1$ 函数；能量有限不意味着速度有界。
\end{example}

\begin{proposition}[分片光滑函数的界面项]\label{prop:weak-interface}
设光滑界面 $\Sigma$ 将区域分成 $\Omega_-,\Omega_+$，法向 $n$ 从负侧指向正侧，$u=u_\pm$ 分片光滑，跳跃 $[u]=u_+-u_-$。其分布导数为
\begin{equation}
 D_i u=(\partial_i u)_{\mathrm{piecewise}}+[u]n_i\,\delta_\Sigma,
 \label{eq:weak-interface}
\end{equation}
这里 $\langle\delta_\Sigma,\varphi\rangle=\int_\Sigma\varphi\dd A$。
\end{proposition}
\begin{proof}
在两侧分别分部积分。负侧界面外法向为 $n$，正侧为 $-n$，故
\[
 -\int u\partial_i\varphi
 =\sum_{\pm}\int_{\Omega_\pm}\partial_i u_\pm\varphi
   +\int_\Sigma(u_+-u_-)n_i\varphi\dd A.
\]
这就是分布作用式。若 $[u]=0$，界面项消失；若有非零跳跃，通常留下无法由体积 $L^p$ 函数表示的曲面测度。
\end{proof}
这解释了连续分片线性有限元为何可属于 $H^1$：函数连续，分片梯度虽有跳跃，但一阶导数没有曲面 Dirac 项。它一般不在 $H^2$，因为梯度的界面跳跃会进入二阶分布导数。

\subsection{一维刻画：弱导数等价于绝对连续代表元}\label{subsec:weak-ac}

\begin{theorem}[一维 Sobolev 微积分基本定理]\label{thm:weak-ac}
在有限区间 $(a,b)$ 上，$u\in W^{1,p}(a,b)$、$1\le p\le\infty$，当且仅当存在代表元 $\widetilde u$ 绝对连续且其 a.e. 经典导数在 $L^p$。此时 $\widetilde u'=Du$ a.e.，并有
\begin{equation}
 \widetilde u(y)-\widetilde u(x)=\int_x^yDu(s)\dd s,
 \qquad x,y\in[a,b].
 \label{eq:weak-ac-ftc}
\end{equation}
\end{theorem}
\begin{proof}
\textbf{从弱导数构造代表元。}令 $g=Du\in L^p\subset L^1$，定义 $v(x)=\int_a^xg(s)\dd s$。Lebesgue 积分的微积分基本定理说明 $v$ 绝对连续，$v'=g$ a.e.。对测试函数用 Fubini：
\[
 \int_a^bv(x)\varphi'(x)\dd x
 =\int_a^bg(s)\left(\int_s^b\varphi'(x)\dd x\right)\dd s
 =-\int_a^bg(s)\varphi(s)\dd s.
\]
所以 $D(u-v)=0$。由命题~\ref{prop:weak-zero-gradient}，$u-v=c$ a.e.，取 $\widetilde u=v+c$ 即得代表元及公式。

\textbf{从绝对连续推出弱导数。}若有绝对连续 $\widetilde u$，则 $\widetilde u\varphi$ 绝对连续，a.e. 导数为 $\widetilde u'\varphi+\widetilde u\varphi'$。积分这个导数，端点项因 $\varphi$ 消失，得到弱导数恒等式。有限区间上的绝对连续函数有界，故本身属于 $L^p$，结论成立。
\end{proof}

当 $p>1$，公式与 H\"older 给出
\[
 |\widetilde u(y)-\widetilde u(x)|
 \le|y-x|^{1-1/p}\|Du\|_p.
\]
$p=1$ 时仍有绝对连续性，但不能从这个估计获得一个正的统一 H\"older 指数。Cantor 函数连续、经典导数 a.e. 为零，却不是绝对连续；由本定理可知它不在 $W^{1,1}(0,1)$。所以“a.e. 存在经典导数且导数可积”仍不足以替代弱导数条件。

在 $(0,1)$ 上 $u=x^\alpha$、$\alpha>0$ 时，直接计算 $u'=\alpha x^{\alpha-1}$，得到 $u\in W^{1,p}\iff\alpha>1-1/p$（$p<\infty$）。等号时出现 $\int_0^1x^{-1}\dd x$。常数 $\alpha=0$ 必须单独处理。

\subsection{沿坐标直线的绝对连续性：高维弱导数的几何含义}

\begin{proposition}[长方体上的切片刻画]\label{prop:weak-acl}
设 $Q=Q'\times(a,b)$，$u,g\in L^p(Q)$，$1\le p<\infty$。则 $D_d u=g$ 当且仅当对几乎所有 $x'\in Q'$，切片 $t\mapsto u(x',t)$ 有绝对连续代表元，导数为 $g(x',t)$ a.e.。
\end{proposition}
\begin{proof}
若 $D_du=g$，用乘积测试函数 $\varphi(x',t)=\eta(x')\psi(t)$，得
\[
 \int_{Q'}\eta(x')\left[\int_a^b
       \big(u(x',t)\psi'(t)+g(x',t)\psi(t)\big)\dd t\right]\dd x'=0.
\]
Fubini 和基本引理说明方括号对 a.e. $x'$ 为零。但这个零测集一开始可能依赖 $\psi$。为统一它，先在每个固定紧支撑区间内取测试函数关于 $C^1$ 范数的可数稠密子集，再对一列穷竭区间取并，去掉对应的可数零测集。其余 $x'$ 上 $u(x',\cdot),g(x',\cdot)\in L^1$，所以该恒等式关于 $\psi$ 的 $C^1$ 范数连续，得以推广到所有测试函数。由一维定理得到绝对连续代表元。

反过来，若 a.e. 切片绝对连续，对任意 $\varphi\in C_c^\infty(Q)$ 先沿 $t$ 做一维分部积分，再用 Fubini 对 $x'$ 积分，便得 $D_du=g$。
\end{proof}
对每个坐标方向重复此论证，得到通常所说的 ACL：沿几乎每条平行坐标轴的线局部绝对连续。这里不要求每条线都成立，也不能据此推出高维函数处处连续。

\subsection{Sobolev 定义、范数与完备性的逐项证明}

\begin{definition}[整数阶 Sobolev 空间]\label{def:sobolev}
对整数 $k\ge0$、$1\le p\le\infty$，定义
\begin{equation}
 W^{k,p}(\Omega)=\{u\in L^p:D^\alpha u\in L^p,\ |\alpha|\le k\},
 \label{eq:sobolev-definition}
\end{equation}
取范数
\begin{equation}
 \|u\|_{W^{k,p}}=
 \begin{cases}(\sum_{|\alpha|\le k}\|D^\alpha u\|_p^p)^{1/p},&p<\infty,\\
 \max_{|\alpha|\le k}\|D^\alpha u\|_\infty,&p=\infty.
 \end{cases}
 \label{eq:sobolev-norm}
\end{equation}
\end{definition}
弱导数的线性直接由定义得到，所以集合对线性运算封闭。各项非负，零阶项保证正定性；齐次性来自积分；三角不等式先对每个导数用 $L^p$ 的 Minkowski，再对有限项向量用 $\ell^p$ 的 Minkowski。只有梯度的表达式通常是半范数，因为它不能识别常数。

记 $H^k=W^{k,2}$，内积为
\begin{equation}
 (u,v)_{H^k}=\sum_{|\alpha|\le k}\int D^\alpha uD^\alpha v\dd x.
 \label{eq:sobolev-inner-product}
\end{equation}
对一般 $p$，采用各梯度分量的 $p$ 次幂之和与采用 $|\nabla u|$ 欧氏长度给出等价范数，常数中可吸收相应维数因子。

\begin{theorem}[Sobolev 完备性]\label{thm:sobolev-complete}
$W^{k,p}(\Omega)$ 是 Banach 空间，$H^k(\Omega)$ 是 Hilbert 空间。
\end{theorem}
\begin{proof}
设 $u_n$ 为 Sobolev Cauchy 序列。对每个 $|\alpha|\le k$，$D^\alpha u_n$ 是 $L^p$ Cauchy，故有极限 $g_\alpha\in L^p$。令 $u=g_0$。关键是确认这些极限之间仍有导数关系。

固定 $\varphi\in C_c^\infty$，由 H\"older，
\[
 \left|\int(u_n-u)D^\alpha\varphi\right|
 \le\|u_n-u\|_p\|D^\alpha\varphi\|_{p'}\to0,
\]
且 $|\int(D^\alpha u_n-g_\alpha)\varphi|\to0$。测试函数紧支撑保证 $p=1$ 所需的 $L^\infty$ 范数、$p=\infty$ 所需的 $L^1$ 范数均有限。在原恒等式中取极限得
\[
 \int uD^\alpha\varphi=(-1)^{|\alpha|}\int g_\alpha\varphi.
\]
所以 $D^\alpha u=g_\alpha$，从而 $u\in W^{k,p}$，有限项收敛又给出 $\|u_n-u\|_{W^{k,p}}\to0$。$p=2$ 时此范数由上述内积诱导，故为 Hilbert 空间。
\end{proof}

同一论证有一个常用的弱极限版本：若 $u_n\to u$ 于分布意义，$D_iu_n\rightharpoonup g_i$ 于 $L^p$（$1<p<\infty$），则 $D_iu=g_i$。测试函数配对即可取极限；不需要导数强收敛。这是 PDE 先验界构造弱解的基本机制。

\subsection{导数与卷积交换：把负号和变量写清楚}

若 $D_i u=g_i$，在 $\operatorname{dist}(x,\partial\Omega)>\varepsilon$ 处，
\begin{align*}
 \partial_i(u*\rho_\varepsilon)(x)
 &=\int u(y)\partial_{x_i}\rho_\varepsilon(x-y)\dd y\\
 &=-\int u(y)\partial_{y_i}\rho_\varepsilon(x-y)\dd y\\
 &=\int g_i(y)\rho_\varepsilon(x-y)\dd y.
\end{align*}
第二行来自核依赖差变量，第三行来自弱导数定义，因此
\begin{equation}
 D_i u_\varepsilon=\rho_\varepsilon*D_i u.
 \label{eq:sobolev-mollifier}
\end{equation}
对 $p<\infty$，近似恒等引理对 $u$ 和每个 $D_i u$ 分别应用，得到 $u_\varepsilon\to u$ 于局部 $W^{1,p}$。

在全空间还可以得到\textbf{紧支撑光滑逼近}：先令 $\chi_R(x)=\chi(x/R)$，$\chi=1$ 于 $B_1$、支撑在 $B_2$，则 $\chi_Ru\to u$ 于 $W^{1,p}(\R^d)$，因为尾部积分趋零，且
\[
 D_i(\chi_Ru)=\chi_RD_iu+(D_i\chi_R)u,
 \quad\|(D_i\chi_R)u\|_p\le C R^{-1}\|u\|_p\to0.
\]
此乘积公式可直接将 $\chi_R\varphi$ 代入定义证明，并不依赖待证的一般乘积定理。再对紧支撑函数卷积，便得 $C_c^\infty(\R^d)$ 在 $W^{1,p}(\R^d)$ 中稠密。

一般开集的内部光滑逼近也可构造：取局部有限光滑单位分解 $\sum_j\psi_j=1$，每个 $\psi_j$ 紧支撑于 $\Omega$。给定 $\eta>0$，逐项将 $\psi_ju$ 卷积，使第 $j$ 项的 $W^{1,p}$ 误差小于 $\eta2^{-j}$，并选择核半径使卷积后的支撑仍局部有限且留在 $\Omega$ 内。局部有限和是光滑函数，总误差不超过 $\eta$。因此 $C^\infty(\Omega)\cap W^{1,p}(\Omega)$ 稠密；通常不能把它改成 $C_c^\infty(\Omega)$。

\begin{warningbox}
上述范数逼近要求 $p<\infty$。$p=\infty$ 时局部卷积仍可用来验证分布恒等式，但不保证 $W^{1,\infty}$ 范数收敛；$|x|$ 的导数跳跃就是反例。
\end{warningbox}

\subsection{乘积法则、链式法则与截断：证明而非形式套用}

\begin{proposition}[有界 Sobolev 函数的乘积]\label{prop:weak-product}
若 $u,v\in W^{1,p}(\Omega)\cap L^\infty(\Omega)$，$1\le p<\infty$，则 $uv\in W^{1,p}$，且 $D_i(uv)=uD_iv+vD_iu$。
\end{proposition}
\begin{proof}
在紧子集内卷积，非负核保证 $|u_\varepsilon|\le\|u\|_\infty$、$|v_\varepsilon|\le\|v\|_\infty$。取一个子列使函数 a.e. 收敛。光滑函数满足通常乘积公式；例如
\[
 u_\varepsilon D_iv_\varepsilon-uD_iv
 =u_\varepsilon(D_iv_\varepsilon-D_iv)
   +(u_\varepsilon-u)D_iv.
\]
第一项以一致上界乘导数的 $L^p$ 收敛趋零；第二项由 a.e. 收敛、上界 $2\|u\|_\infty|D_iv|$ 和控制收敛在 $L^p$ 中趋零。另一项同理。对测试函数取极限，得到弱乘积公式；右侧全局属于 $L^p$，而 $uv\in L^p$，故是全局结论。
\end{proof}

\begin{proposition}[有界导数的链式法则]\label{prop:weak-chain}
若 $F\in C^1(\R)$、$\|F'\|_\infty<\infty$，$u\in W^{1,p}(\Omega)$、$p<\infty$，则局部有 $D_iF(u)=F'(u)D_i u$。若 $F(0)=0$ 或 $|\Omega|<\infty$，则 $F(u)\in W^{1,p}(\Omega)$。
\end{proposition}
\begin{proof}
局部光滑逼近后，$F(u_\varepsilon)\to F(u)$ 于 $L^p$，因为 $F$ 是 Lipschitz。对导数作分解
\[
 F'(u_\varepsilon)D_i u_\varepsilon-F'(u)D_i u
 =F'(u_\varepsilon)(D_i u_\varepsilon-D_i u)
  +(F'(u_\varepsilon)-F'(u))D_i u.
\]
第一项强收敛，第二项取 a.e. 收敛子列后由 $F'$ 连续及控制收敛得到 $L^p$ 收敛。因此积分恒等式可取极限。最后 $|F(u)|\le|F(0)|+\|F'\|_\infty|u|$，给出全局可积性条件。
\end{proof}

\textbf{绝对值与正部。}先取 $F_\varepsilon(s)=\sqrt{s^2+\varepsilon^2}-\varepsilon$。其导数有界于 1，$F_\varepsilon(u)\to|u|$ 于 $L^p$，且 $F_\varepsilon'(u)D_i u\to\operatorname{sgn}(u)D_i u$ 于 $L^p$，其中 $\operatorname{sgn}(0)=0$。所以 $D_i|u|=\operatorname{sgn}(u)D_i u$。

还需要说明 $D_i u=0$ a.e. 于 $\{u=0\}$。由切片刻画，将问题归为一维绝对连续函数：在零点集的几乎每个点，该点既是零点集的密度点，又是可微点；取零点序列趋向它，差商为零，所以导数为零。Lebesgue 密度定理和 Fubini 给出结论。于是
\[
 D_i u^+=\mathbf1_{\{u>0\}}D_i u,\qquad
 D_iT_M(u)=\mathbf1_{\{|u|<M\}}D_i u,
\]
其中 $u^+=\max(u,0)$、$T_M(s)=\max(-M,\min(s,M))$。截断公式也可由绝对连续切片直接验证。它们用于最大值原理、非负性保持与弱解测试。

\subsection{差商判据：不用先知道导数也能检测 Sobolev 正则性}

令 $\delta_h^iu(x)=[u(x+he_i)-u(x)]/h$。若线段落在定义域内，光滑函数的基本定理给出
\[
 \delta_h^iu(x)=\int_0^1D_i u(x+the_i)\dd t.
\]
由 Jensen 积分并作变量替换，在全空间有
\begin{equation}
 \|\delta_h^iu\|_p\le\|D_i u\|_p,
 \qquad u\in W^{1,p}(\R^d),\quad1\le p<\infty.
 \label{eq:weak-difference-bound}
\end{equation}
对一般 Sobolev 函数由光滑稠密性推广。局部版本用包含所有平移线段的较大紧子域控制。

\begin{theorem}[局部差商的反向判据]\label{thm:weak-difference}
设 $1<p<\infty$，$u\in L^p_{\mathrm{loc}}(\Omega)$。若对每个 $U\Subset\Omega$、每个方向 $i$，存在 $C_U$，使所有充分小 $h$ 满足 $\|\delta_h^iu\|_{L^p(U)}\le C_U$，则 $u\in W^{1,p}_{\mathrm{loc}}(\Omega)$。
\end{theorem}
\begin{proof}
固定 $U$，由 $L^p$ 自反性，任取 $h_n\to0$ 可抽取子列，使 $\delta_{h_n}^iu\rightharpoonup g_i$ 于 $L^p(U)$。对 $\varphi\in C_c^\infty(U)$，变量替换得到
\[
 \int\delta_h^iu\,\varphi
 =-\int u\,\delta_{-h}^i\varphi
 \longrightarrow-\int u\partial_i\varphi.
\]
左侧同时趋于 $\int g_i\varphi$，故 $g_i=D_i u$。在穷竭子域上逐一构造，唯一性使这些局部导数相容。弱下半连续还给出 $\|D_i u\|_{L^p(U)}\le C_U$。
\end{proof}

$p=1$ 的反向判据不能直接成立：阶跃函数的差商 $L^1$ 范数一致有界，却没有 $L^1$ 弱导数；它对应 BV（有界变差）空间中的测度导数。差商方法也会用于椭圆方程的二阶弱正则性估计。

% 接在弱导数理论之后；边界构造的证明限定在有界 C^1 区域。
\subsection{为什么需要延拓：先在平边界上亲手构造}

对 $\R^d_+=\{x_d>0\}$，定义偶反射
\[
 Eu(x',x_d)=u(x',|x_d|).
\]
若 $u\in W^{1,p}(\R^d_+)$、$1\le p<\infty$，由切片刻画，一维法向切片在有限区间上有绝对连续代表元，偶反射后仍绝对连续。切向切片只是在另一侧复制。因此反射函数的弱导数为
\begin{equation}
 D_jEu(x',x_d)=
 \begin{cases}(D_ju)(x',|x_d|),&j<d,\\
 \operatorname{sgn}(x_d)(D_du)(x',|x_d|),&j=d.
 \end{cases}
 \label{eq:weak-reflection}
\end{equation}
切片上的导数与弱导数等价，故这不是仅对光滑函数的形式计算。上下半空间作变量替换得到
\[
 \|Eu\|_p^p=2\|u\|_p^p,\qquad
 \|D_jEu\|_p^p=2\|D_ju\|_p^p.
\]
因此反射是有界线性延拓。注意这里不能使用零延拓代替：若上侧边界值非零，零延拓产生界面跳跃与 Dirac 曲面项。

\subsection{有界 \texorpdfstring{$C^1$}{C1} 区域的延拓：坐标拉平与单位分解}\label{subsec:weak-extension}

\begin{theorem}[一阶有界延拓]\label{thm:weak-extension}
若 $\Omega\subset\R^d$ 为有界 $C^1$ 区域，$1\le p<\infty$，存在有界线性映射
\[
 E:W^{1,p}(\Omega)\longrightarrow W^{1,p}(\R^d),\qquad
 Eu=u\text{ a.e. 于 }\Omega,
\]
且可使所有 $Eu$ 的支撑包含于同一个有界集内。
\end{theorem}
\begin{proof}
\textbf{第一步：写出坐标变换。}边界紧，可用有限个坐标邻域覆盖。经刚体变换，在每个邻域内
$\Omega=\{(x',x_d):x_d>\gamma(x')\}$，$\gamma\in C^1$。定义
\[
 \Phi(y',s)=(y',s+\gamma(y')),\qquad
 \Phi^{-1}(x',x_d)=(x',x_d-\gamma(x')).
\]
两者在选定的小邻域内导数有界，Jacobian 行列式为 1。令 $v=u\circ\Phi$，链式变换为
\[
 \partial_s v=(D_du)\circ\Phi,\qquad
 \partial_{y_j}v=(D_ju)\circ\Phi+
                 (\partial_j\gamma)(D_du)\circ\Phi\quad(j<d).
\]
这些公式对弱导数也成立：在任何紧含于上半邻域的集合中先光滑逼近 $u$，对光滑函数计算，再用变量替换和有界 Jacobian 将函数及一阶导数的 $L^p$ 收敛传递过去，最后在测试函数恒等式中取极限。穷竭紧子集便覆盖上半邻域。故坐标拉平及拉回均有 $W^{1,p}$ 有界估计。

\textbf{第二步：对局部块反射。}取一组光滑单位分解 $\chi_0,\ldots,\chi_N$，在 $\overline\Omega$ 的邻域内和为 1；$\chi_0$ 支撑远离边界，其余各项支撑在上述坐标邻域内。有限光滑单位分解是这里使用的标准几何工具。对 $\chi_j u$ 拉平，在平边界上偶反射，再拉回；选择邻域与截断使该局部块在坐标邻域外缘附近为零，然后可以向外零延拓。利用上一小节反射估计、坐标变换和光滑截断乘积公式，得到
\[
 \|E_j(\chi_j u)\|_{W^{1,p}(\R^d)}
 \le C_j\|u\|_{W^{1,p}(\Omega\cap U_j)}.
\]
$\chi_0u$ 直接零延拓即可，因为支撑远离边界，不会有界面项。

\textbf{第三步：求和。}将这些有限个线性局部延拓相加，便得 $Eu$。在 $\Omega$ 内和为 $u\sum_j\chi_j=u$，且总范数由有限个常数之和控制。所有截断、坐标邻域预先固定且有界，所以延拓的支撑也统一有界。
\end{proof}

这也证明了 $C^\infty(\overline\Omega)$ 在 $W^{1,p}(\Omega)$ 中稠密：将 $Eu$ 在全空间卷积，再限制回 $\Omega$，即可在函数和全部一阶导数的 $L^p$ 范数下收敛。这里 $C^\infty(\overline\Omega)$ 指在闭包邻域内光滑的函数的限制，允许边界值非零。

延拓结论也适用于 Lipschitz 区域，但图函数只有 a.e. 导数时，须补上 Lipschitz 变换的弱链式法则。为保持本节每一步可验证，上面的完整构造限定为 $C^1$ 边界；本专题后续需要延拓的完整证明均采用该假设。

\subsection{迹不等式与迹算子：先证明平边界估计}

取柱体 $Q=Q'\times(0,L)$，先设 $u$ 在闭柱体附近光滑。对 $p>1$，沿竖直线由基本定理得
\[
 |u(x',0)|^p
 \le |u(x',t)|^p+
 p\int_0^L|u(x',s)|^{p-1}|\partial_du(x',s)|\dd s.
\]
对 $t$ 在 $(0,L)$ 上取平均，再对 $x'$ 积分，H\"older 给出
\begin{align*}
 \|u(\cdot,0)\|_{L^p(Q')}^p
 &\le L^{-1}\|u\|_{L^p(Q)}^p
    +p\|u\|_{L^p(Q)}^{p-1}\|\partial_du\|_{L^p(Q)}\\
 &\le C_{p,L}\big(\|u\|_{L^p(Q)}^p+\|\partial_du\|_{L^p(Q)}^p\big).
\end{align*}
最后一步是 Young。当 $p=1$，直接使用
$|u(x',0)|\le|u(x',t)|+\int_0^L|\partial_du|$，得到同样的范数控制。拉平边界后，曲面测度为 $\sqrt{1+|\nabla\gamma|^2}\dd x'$，因子有界，结合有限单位分解得到一般边界估计。

\begin{theorem}[迹算子的构造]\label{thm:weak-trace}
若 $\Omega$ 有界 $C^1$，$1\le p<\infty$，存在唯一有界线性算子
\begin{equation}
 \operatorname{Tr}:W^{1,p}(\Omega)\to L^p(\partial\Omega),\qquad
 \|\operatorname{Tr}u\|_{L^p(\partial\Omega)}\le C\|u\|_{W^{1,p}},
 \label{eq:weak-trace-estimate}
\end{equation}
它对光滑函数等于通常边界限制。
\end{theorem}
\begin{proof}
取光滑逼近 $u_n\to u$ 于 $W^{1,p}$。刚才的估计用于差值，说明 $u_n|_{\partial\Omega}$ 在 $L^p(\partial\Omega)$ 中 Cauchy。用 $L^p$ 完备性定义其极限为 $\operatorname{Tr}u$。若换一组逼近，同一估计用于两组之差，说明极限无关。线性与估计取极限保留，光滑稠密性保证这样的连续延伸唯一。

还可确认这个迹等于 a.e. 法向绝对连续切片的端点值：在平柱体上，选取误差的 $W^{1,p}$ 范数的 $p$ 次幂可求和的光滑子列。Fubini 说明对 a.e. $x'$，其切片在一维 $W^{1,p}(0,L)$ 中收敛；一维端点估计又保证端点值收敛到该绝对连续切片的端点值。这与已经定义的 $L^p$ 边界极限相同。
\end{proof}
边界在体积测度下是零测集，所以迹不能通过“挑一个代表元直接在边界取值”定义。对于 $p=2$，更精确的边界目标空间是 $H^{1/2}(\partial\Omega)$；这一分数阶迹定理此处仅作延伸方向，不作证明，也不宣称任意 $L^2$ 边界数据都有 $H^1$ 延拓。

\subsection{零边界空间为什么恰好是零迹空间}

定义
\begin{equation}
 W^{1,p}_0(\Omega)=\overline{C_c^\infty(\Omega)}^{W^{1,p}},
 \qquad H^1_0=W^{1,2}_0.
 \label{eq:sobolev-zero}
\end{equation}
它作为完备空间的闭子空间也完备。对有界 $C^1$ 区域和 $p<\infty$，有
\begin{equation}
 \boxed{W^{1,p}_0(\Omega)=\ker\operatorname{Tr}.}
 \label{eq:weak-trace-kernel}
\end{equation}

\begin{proof}
\textbf{正向包含。}紧支撑光滑函数的迹为零，由迹算子连续性，其 Sobolev 闭包的迹也为零。

\textbf{反向包含先看平边界。}若上半柱体的 $u$ 迹为零，且已作截断使其在侧面和顶面附近为零，那么几乎每条法向绝对连续切片的下端值为零。将切片向下零延拓，仍是绝对连续函数，导数是原导数的零延拓；切向导数由分部积分/Fubini 保留。因此全空间零延拓 $Zu$ 属于 $W^{1,p}(\R^d)$，没有边界测度项。

令 $w_h(y',s)=Zu(y',s-h)$，$h>0$，即把支撑向域内移动 $h$。全空间平移连续性对函数及导数同时成立，故 $w_h\to Zu$ 于 $W^{1,p}$。再以半径小于 $h/2$ 的核卷积，得到支撑距平边界为正的光滑逼近。

\textbf{再处理曲边界。}对零迹 $u$ 作前述有限单位分解并拉平边界。每个边界块按上述方法零延拓并内移，再拉回，得到支撑紧含于原区域的 $W^{1,p}$ 逼近。$C^1$ 坐标拉回不一定保持无穷光滑，但支撑已远离边界，此时再在原坐标内卷积即可得到真正的 $C_c^\infty(\Omega)$ 逼近。内域块直接卷积。有限求和并逐项控制误差得 $u\in W^{1,p}_0$。
\end{proof}

这个论证还说明：$u\in W^{1,p}_0$ 的零延拓属于全空间 $W^{1,p}$，并且 $D_iZu=Z(D_i u)$。对任意开集，这个结论也可直接用 $C_c^\infty$ 逼近的定义证明。无滑移边界取 $\vct u\in H^1_0(\Omega)^d$，正是同时要求有限能量和零迹。

\subsection{弱紧性、下半连续性及本节的依赖边界}

$W^{k,p}$ 通过 $u\mapsto(D^\alpha u)_{|\alpha|\le k}$ 等距嵌入有限个 $L^p$ 的乘积，像由完备性证明可知是闭的。因此 $p<\infty$ 时可分，$1<p<\infty$ 时自反；这里使用了 $L^p$ 的相应结构定理。故有界 Sobolev 序列在 $1<p<\infty$ 时可抽出弱收敛子列。

若 $u_n\rightharpoonup u$ 于赋范空间，利用 Hahn--Banach 给出的对偶范数表示，可选 $\ell$ 满足 $\|\ell\|=1$、$\ell(u)=\|u\|$，于是
\[
 \|u\|=\lim_n\ell(u_n)\le\liminf_n\|u_n\|.
\]
这就是范数的弱下半连续性，能量泛函由此在弱极限下得到下界。

本节完整展开了一阶弱导数、卷积、切片、运算、差商及 $C^1$ 边界上的一阶延拓与迹构造。作为预备引用的工具包括 Lebesgue 积分与微分基础、测度正则性、光滑单位分解、$L^p$ 对偶/自反性和 Hahn--Banach；没有把这些基础定理也当作已经在本节重新证明。



% 从一阶基础不等式到临界、Morrey 和 Rellich；保留旧交叉引用。
\Needspace{12\baselineskip}
\section{Sobolev 嵌入：次临界、临界、连续性与紧性的详细证明}\label{sec:sobolev-embedding}

\subsection{先区分包含、连续嵌入与紧嵌入}

$X\subset Y$ 只说明每个 $X$ 元素也属于 $Y$。$X\hookrightarrow Y$ 还要求恒等映射连续，即存在与 $u$ 无关的 $C$，使 $\|u\|_Y\le C\|u\|_X$。这保证 $X$ 中强收敛传递到 $Y$。$X\Subset Y$ 表示紧嵌入：$X$ 中每个有界序列都有 $Y$ 中的强收敛子列。

接下来的证明按四条路线展开：$p<d$ 得更高可积性，$p=d$ 得任意有限次幂可积性，$p>d$ 得 H\"older 连续代表元，最后以平移估计证明紧性。需要一般边界延拓时，完整证明统一采用上一节已经构造过的有界 $C^1$ 区域。

\begin{theorem}[次临界一阶 Sobolev 嵌入]\label{thm:sobolev-main}
设整数 $d\ge2$，$1\le p<d$，$p^*=dp/(d-p)$。对任意开集 $\Omega\subset\R^d$，
\begin{equation}
 \boxed{\|u\|_{p^*}\le C_{d,p}\|\nabla u\|_p,
       \qquad u\in W^{1,p}_0(\Omega).}
 \label{eq:sobolev-main}
\end{equation}
在全空间 $W^{1,p}(\R^d)$ 上也有同一估计。若 $\Omega$ 有界且为 $C^1$，一般非零边界函数满足
\begin{equation}
 \|u\|_{p^*}\le C_{\Omega,p}\big(\|u\|_p+\|\nabla u\|_p\big),
 \qquad u\in W^{1,p}(\Omega).
 \label{eq:sobolev-nonzero}
\end{equation}
\end{theorem}

\subsection{基础积分引理：将不同坐标缺失的函数乘起来}

\begin{lemma}[坐标乘积积分估计]\label{lem:sobolev-coordinate}
对每个 $i=1,\ldots,d$，设非负可积函数 $A_i(\widehat x_i)$ 只依赖除 $x_i$ 以外的 $d-1$ 个变量，则
\begin{equation}
 \int_{\R^d}\prod_{i=1}^d A_i(\widehat x_i)^{1/(d-1)}\dd x
 \le\prod_{i=1}^d
       \left(\int_{\R^{d-1}}A_i\dd\widehat x_i\right)^{1/(d-1)}.
 \label{eq:sobolev-coordinate}
\end{equation}
\end{lemma}
\begin{proof}
对维数归纳。$d=2$ 时是 $\iint A_1(x_2)A_2(x_1)=\int A_1\int A_2$。

设 $d\ge3$，并假设 $d-1$ 维成立。先沿 $x_d$ 积分，$A_d$ 与它无关。对其余 $d-1$ 个因子用指数全为 $d-1$ 的 H\"older，令
\[
 G_i(\widehat x_i')=\int_\R A_i(\widehat x_i)\dd x_d
 \quad(1\le i<d),\qquad x'=(x_1,\ldots,x_{d-1}).
\]
则原积分不超过
\[
 \int_{\R^{d-1}} A_d(x')^{1/(d-1)}
                \prod_{i<d}G_i(\widehat x_i')^{1/(d-1)}\dd x'.
\]
设 $H(x')=\prod_{i<d}G_i^{1/(d-2)}$，上式的后半乘积为 $H^{(d-2)/(d-1)}$。再用共轭指数 $d-1$、$(d-1)/(d-2)$ 的 H\"older，上式不超过
\[
 \left(\int A_d\right)^{1/(d-1)}
 \left(\int H\right)^{(d-2)/(d-1)}.
\]
对 $H$ 使用 $d-1$ 维归纳假设，即 $\int H\le\prod_{i<d}(\int G_i)^{1/(d-2)}$。而 Tonelli 给出 $\int G_i=\int A_i$。代入即可得到所需的全部 $d$ 个因子。
\end{proof}

\subsection{第一步：从微积分基本定理证明 \texorpdfstring{$W^{1,1}$}{W11} 估计}

设 $v\in C_c^1(\R^d)$。固定除 $x_i$ 以外的变量，从负无穷积分，因紧支撑有
\[
 |v(x)|=\left|\int_{-\infty}^{x_i}\partial_i v(x_1,\ldots,s,\ldots,x_d)\dd s\right|
 \le A_i(\widehat x_i):=\int_\R|\partial_i v|\dd x_i.
\]
将 $d$ 个不等式相乘并取 $1/(d-1)$ 次幂，得
\[
 |v(x)|^{d/(d-1)}\le\prod_{i=1}^d A_i(\widehat x_i)^{1/(d-1)}.
\]
积分并用刚证明的引理，再取 $(d-1)/d$ 次幂：
\begin{align*}
 \|v\|_{d/(d-1)}
 &\le\prod_{i=1}^d\|\partial_i v\|_1^{1/d}\\
 &\le\frac1d\sum_{i=1}^d\|\partial_i v\|_1
 \le C_d\|\nabla v\|_1.
\end{align*}
于是
\begin{equation}
 \boxed{\|v\|_{d/(d-1)}\le C_d\|\nabla v\|_1.}
 \label{eq:sobolev-w11}
\end{equation}
这是整个次临界嵌入证明的起点，且已经覆盖 $p=1$。

\textbf{二维单独核对。}此时 $|v(x,y)|$ 同时不超过 $A(y)$ 与 $B(x)$，所以
\[
 |v|^2\le A(y)B(x),\qquad
 \|v\|_2^2\le\|\partial_1v\|_1\|\partial_2v\|_1.
\]

\textbf{三维单独核对。}令 $A(y,z)=\int|\partial_1v|\dd x$、$B(x,z)=\int|\partial_2v|\dd y$、$C(x,y)=\int|\partial_3v|\dd z$。则
\[
 \iiint|v|^{3/2}\le\iiint(ABC)^{1/2}
 \le(\|\partial_1v\|_1\|\partial_2v\|_1\|\partial_3v\|_1)^{1/2}.
\]
最后一个估计就是引理在 $d=3$ 时的三次 Cauchy--Schwarz。取 $2/3$ 次幂得到 $\|v\|_{3/2}\le C\|\nabla v\|_1$。

\subsection{第二步：为什么必须选这个幂指数}

设 $1<p<d$，先取 $u\in C_c^\infty(\R^d)$。希望把基础估计应用于 $v=|u|^\gamma$。左边会出现 $\|u\|_{\gamma d/(d-1)}^\gamma$；右边的导数乘积在 H\"older 后会出现 $\|u\|_{(\gamma-1)p'}^{\gamma-1}$。为了两侧使用同一个函数范数，要求
\[
 \frac{\gamma d}{d-1}=(\gamma-1)p'.
\]
这不是猜测，而是解一个代数方程。代入 $p'=p/(p-1)$，交叉相乘得
\[
 \gamma d(p-1)=p(d-1)(\gamma-1),\qquad
 \gamma(d-p)=p(d-1).
\]
所以
\begin{equation}
 \gamma=\frac{p(d-1)}{d-p}>1,\qquad
 \frac{\gamma d}{d-1}=(\gamma-1)p'=p^*=\frac{dp}{d-p}.
 \label{eq:sobolev-power-choice}
\end{equation}
因为 $\gamma>1$，$|s|^\gamma$ 在零点也 $C^1$，故 $v\in C_c^1$；不需要宣称它无穷可微。逐项计算得
\begin{align*}
 \|u\|_{p^*}^{\gamma}
 &\le C_d\|\nabla(|u|^\gamma)\|_1\\
 &=C_d\gamma\int|u|^{\gamma-1}|\nabla u|\dd x\\
 &\le C_d\gamma\|\nabla u\|_p
      \left(\int|u|^{(\gamma-1)p'}\dd x\right)^{1/p'}\\
 &=C_d\gamma\|\nabla u\|_p\|u\|_{p^*}^{\gamma-1}.
\end{align*}
若 $u$ 非零可除去 $\|u\|_{p^*}^{\gamma-1}$；若为零结论直接成立。这就证明了光滑函数版本。三维 $p=2$ 对应 $\gamma=4$：
\[
 \|u\|_6^4\le C\int|u|^3|\nabla u|
 \le C\|u\|_6^3\|\nabla u\|_2.
\]

\subsection{第三步：光滑函数的估计怎样推广到 Sobolev 函数}

\textbf{零边界情形。}若 $u\in W^{1,p}_0(\Omega)$，按定义取 $u_n\in C_c^\infty(\Omega)$，使其在 $W^{1,p}$ 范数中趋于 $u$。将各 $u_n$ 零延拓，已证估计用于差值：
\[
 \|u_n-u_m\|_{p^*}\le C\|\nabla u_n-\nabla u_m\|_p\to0.
\]
由 $L^{p^*}$ 完备性得到某个 $w\in L^{p^*}$。还须确认 $w=u$：对任意有限测度 $K\Subset\Omega$，
\[
 \|u-w\|_{L^1(K)}
 \le |K|^{1-1/p}\|u-u_n\|_{L^p(K)}
     +|K|^{1-1/p^*}\|u_n-w\|_{L^{p^*}(K)}\to0.
\]
故 $u=w$ a.e.，用可数穷竭覆盖区域即可。最后范数收敛允许在不等式两侧取极限，证明~\eqref{eq:sobolev-main}。

\textbf{全空间情形。}上一节已经构造了 $C_c^\infty(\R^d)$ 的 $W^{1,p}$ 逼近，因此逐字重复上面的 Cauchy 与极限识别论证即可。

\textbf{一般边界情形。}对有界 $C^1$ 区域，使用定理~\ref{thm:weak-extension} 的延拓 $Eu$，则
\[
 \|u\|_{L^{p^*}(\Omega)}
 \le\|Eu\|_{L^{p^*}(\R^d)}
 \le C\|\nabla Eu\|_{L^p(\R^d)}
 \le C_\Omega\|u\|_{W^{1,p}(\Omega)}.
\]
延拓中的截断产生零阶项，所以不能删掉 $\|u\|_p$。非零常数也是立即可见的反例。至此定理~\ref{thm:sobolev-main} 的三个版本全部证明。

\subsection{缩放推导临界指数，并构造不能改进的反例}

令 $u_\lambda(x)=u(\lambda x)$，变量替换给出
\[
 \|u_\lambda\|_q=\lambda^{-d/q}\|u\|_q,\qquad
 \|\nabla u_\lambda\|_p=\lambda^{1-d/p}\|\nabla u\|_p.
\]
若全空间的纯梯度估计 $\|u\|_q\le C\|\nabla u\|_p$ 对所有尺度成立且 $C$ 不变，则必须 $-d/q=1-d/p$，即 $q=p^*$。

在有界区域内部一点 $x_0$ 附近，取非零 $\phi\in C_c^\infty(B_1)$，定义
\[
 u_n(x)=n^{d/p-1}\phi(n(x-x_0)),
\]
对充分大的 $n$ 支撑落在区域内。计算
\begin{equation}
 \|\nabla u_n\|_p=\|\nabla\phi\|_p,\quad
 \|u_n\|_p=n^{-1}\|\phi\|_p,\quad
 \|u_n\|_q=n^{d/p-1-d/q}\|\phi\|_q.
 \label{eq:sobolev-concentration}
\end{equation}
若 $q>p^*$，最后指数为正，$L^q$ 范数发散，所以这种更高指数的连续嵌入一般不成立。若 $q=p^*$，范数恒为正，而 $u_n\to0$ a.e.，故不可能有 $L^{p^*}$ 强收敛子列。这同时解释端点连续但一般不紧。

\subsection{临界情况 \texorpdfstring{$p=d$}{p=d}：有限次幂可以，端点无穷不可以}

\begin{theorem}[临界有限指数嵌入]\label{thm:sobolev-critical}
若 $d\ge2$，$\Omega$ 有界 $C^1$，则对每个 $1\le q<\infty$，
\[
 \|u\|_q\le C_{\Omega,q}\|u\|_{W^{1,d}},
 \qquad u\in W^{1,d}(\Omega).
\]
\end{theorem}
\begin{proof}
$q\le d$ 时由有限测度上的 $L^d\hookrightarrow L^q$ 直接成立。若 $q>d$，选
\[
 r=\frac{dq}{d+q},\qquad1<r<d,\qquad\frac{dr}{d-r}=q.
\]
有限测度上的 H\"older 给出 $W^{1,d}\hookrightarrow W^{1,r}$。对 $r$ 使用已证的次临界一般边界嵌入：
\[
 \|u\|_q\le C_{\Omega,r}\|u\|_{W^{1,r}}
 \le C_{\Omega,q}\|u\|_{W^{1,d}}.
\]
这说明每个有限 $q$ 都可以，但常数依赖 $q$，没有推出 $q=\infty$ 的估计。
\end{proof}

\begin{example}[临界无界函数的完整检查]\label{ex:sobolev-loglog}
在 $d\ge2$、原点附近取 $u(x)=\log\log(e/|x|)$，再乘支撑于 $B_1$、原点附近为 1 的光滑截断。令 $r=|x|$，有
\[
 |u'(r)|=\frac1{r\log(e/r)},\qquad
 \int_0^\delta|u'(r)|^dr^{d-1}\dd r
 =\int_0^\delta\frac{\dd r}{r\log^d(e/r)}<\infty.
\]
最后的积分以 $s=\log(e/r)$ 化为 $\int^\infty s^{-d}\dd s$，因 $d>1$ 收敛。函数本身的 $d$ 次幂积分也收敛：同一换元后被指数衰减 $e^{-ds}$ 控制。去掉 $B_\varepsilon$ 分部积分时，内边界项不超过 $C\varepsilon^{d-1}\log\log(e/\varepsilon)\to0$，所以候选导数确实是弱导数，无额外原点质量。因此 $u\in W^{1,d}_0(B_1)$，但 $u$ 本质无界。特别二维 $H^1$ 不嵌入 $L^\infty$。
\end{example}

\subsection{Morrey 证明的核心：从线段平均得到奇核估计}

设 $u\in C_c^\infty(\R^d)$，$B_r(x)$ 为球，记 $u_B=|B|^{-1}\int_Bu$。从沿线段的基本定理出发，
\[
 |u(x)-u(x+s\omega)|\le\int_0^s|\nabla u(x+t\omega)|\dd t.
\]
在球上取平均，极坐标和 Tonelli 给出
\begin{align*}
 |u(x)-u_{B_r(x)}|
 &\le\frac1{|B_r|}\int_{\mathbb S^{d-1}}\int_0^r
       \int_0^s|\nabla u(x+t\omega)|\dd t\,s^{d-1}\dd s\dd\omega\\
 &\le C_d\int_{\mathbb S^{d-1}}\int_0^r
                   |\nabla u(x+t\omega)|\dd t\dd\omega\\
 &=C_d\int_{B_r(x)}\frac{|\nabla u(z)|}{|x-z|^{d-1}}\dd z.
\end{align*}
第二行使用 $|B_r|^{-1}\int_t^rs^{d-1}\dd s\le C_d$。这个核的分母次数来自极坐标体积元 $t^{d-1}$，不能凭空省略。

若 $p>d$，则 $p'=p/(p-1)$ 满足 $(d-1)p'<d$。H\"older 的核积分为
\begin{align*}
 \left(\int_{B_r(x)}|x-z|^{-(d-1)p'}\dd z\right)^{1/p'}
 &=\left(|\mathbb S^{d-1}|\int_0^r
                 t^{d-1-(d-1)p'}\dd t\right)^{1/p'}\\
 &=C_{d,p}r^{1-d/p}.
\end{align*}
因此
\begin{equation}
 |u(x)-u_{B_r(x)}|\le C_{d,p}r^{1-d/p}\|\nabla u\|_{L^p(B_r(x))}.
 \label{eq:sobolev-morrey-average}
\end{equation}
这里恰好需要 $p>d$；当 $p=d$ 核积分临界发散，这说明证明在哪里失效。

\subsection{Morrey 连续嵌入：两点估计、上界与代表元}

\begin{theorem}[Morrey 不等式]\label{thm:sobolev-morrey}
若 $d<p<\infty$，$\alpha=1-d/p$，则 $W^{1,p}(\R^d)$ 的每个 a.e. 等价类有唯一连续代表元，满足
\begin{equation}
 [u]_{C^{0,\alpha}}:=\sup_{x\ne y}\frac{|u(x)-u(y)|}{|x-y|^\alpha}
 \le C\|\nabla u\|_p,\qquad
 \|u\|_\infty\le C(\|u\|_p+\|\nabla u\|_p).
 \label{eq:sobolev-morrey}
\end{equation}
有界 $C^1$ 区域上有 $W^{1,p}(\Omega)\hookrightarrow C^{0,\alpha}(\overline\Omega)$，估计由完整 Sobolev 范数控制。
\end{theorem}
\begin{proof}
\textbf{两点之差。}先设 $u$ 光滑。取 $x\ne y$、$r=2|x-y|$，令 $E=B_r(x)\cap B_r(y)$。交集包含半径 $3r/4$、中心为 $(x+y)/2$ 的球，因此 $|E|\ge c_dr^d$。对 $z\in E$ 写
\[
 |u(x)-u(y)|\le |u(x)-u(z)|+|u(y)-u(z)|,
\]
在 $E$ 上平均。第一项不超过 $C_dr^{-d}\int_{B_r(x)}|u(x)-u(z)|\dd z$，第二项类似。沿上一小节的线段平均计算再用 H\"older，即得
$|u(x)-u(y)|\le C r^\alpha\|\nabla u\|_p=C'|x-y|^\alpha\|\nabla u\|_p$。

\textbf{本身的上界。}取半径 1 的球，利用~\eqref{eq:sobolev-morrey-average} 及 $|u_{B_1(x)}|\le |B_1|^{-1/p}\|u\|_p$，得到一致的 $\|u\|_\infty$ 估计。

\textbf{推广与代表元。}取 $u_n\in C_c^\infty$ 使其在 $W^{1,p}$ 中趋于 $u$。上述两个估计用于 $u_n-u_m$，说明它们在完整 $C^{0,\alpha}$ 范数下 Cauchy，特别一致收敛到某个连续函数 $\widetilde u$。它在局部 $L^p$ 的极限与 $u$ 一致，所以是代表元；不等式取极限保留。两个连续代表元若 a.e. 相同，差若某点非零则在一个球内非零，与 a.e. 相同矛盾，故代表元唯一。一般边界情形先延拓到全空间，再限制至闭包即可。
\end{proof}

当 $p=\infty$，全空间 $W^{1,\infty}$ 具有 Lipschitz 代表元。可在光滑卷积上用线段基本定理得到 $|u_\varepsilon(x)-u_\varepsilon(y)|\le\|\nabla u\|_\infty|x-y|$，再用局部一致紧性与局部 $L^1$ 收敛识别代表元；这里无需 $W^{1,\infty}$ 范数的光滑稠密性。凸区域内也能直接沿线段估计；一般区域应考虑延拓或域的几何性质。

\subsection{Poincar\'e：零边界版本的独立证明}

\begin{theorem}[零边界 Poincar\'e]\label{thm:sobolev-poincare}
若 $\Omega$ 有界，$1\le p<\infty$，则
\begin{equation}
 \|u\|_p\le C_\Omega\|\nabla u\|_p,
 \qquad u\in W^{1,p}_0(\Omega).
 \label{eq:sobolev-poincare}
\end{equation}
\end{theorem}
\begin{proof}
将区域平移后放进边长 $L$ 的立方体 $Q=(0,L)^d$。对 $u\in C_c^\infty(\Omega)$ 零延拓，有
\[
 u(x_1,x')=\int_0^{x_1}\partial_1u(s,x')\dd s.
\]
H\"older（$p=1$ 直接积分）给出
\[
 |u(x_1,x')|^p\le L^{p-1}\int_0^L|\partial_1u(s,x')|^p\dd s.
\]
再积分 $x_1,x'$ 得 $\|u\|_p^p\le L^p\|\partial_1u\|_p^p$。按 $W^{1,p}_0$ 定义逼近并取极限即可。这个证明不需要边界光滑，也不依赖后面的紧性定理。
\end{proof}

因此 $\|\nabla u\|_p$ 在零边界空间上是等价范数；在一般 $W^{1,p}$ 上仍只是半范数。三维的 Sobolev 嵌入再结合有限测度的 $L^6\hookrightarrow L^2$ 也能推出 $p=2$ 的零边界 Poincar\'e，但沿线积分证明适用范围更直接。

% 延续嵌入章节：空间紧性、零均值 Poincare 和高阶补充。
\subsection{Rellich 的第一步：一致平移估计与一致光滑逼近}

对 $f\in W^{1,p}(\R^d)$，$1\le p<\infty$，沿向量 $h$ 的线段积分并用 Jensen，得到
\begin{equation}
 \|f(\cdot-h)-f\|_p\le |h|\|\nabla f\|_p.
 \label{eq:sobolev-translation-uniform}
\end{equation}
光滑时逐点有 $f(x-h)-f(x)=-\int_0^1\nabla f(x-th)\cdot h\dd t$，积分即可；一般情况用光滑稠密性推广。将此式用于卷积平均，得到定量估计
\begin{equation}
 \|f*\rho_\varepsilon-f\|_p
 \le C_\rho\varepsilon\|\nabla f\|_p.
 \label{eq:sobolev-uniform-mollify}
\end{equation}
与普通近似恒等不同，这个误差对\textbf{整个梯度范数有界的函数族一致趋零}，它是紧性证明的核心。

设 $u_n$ 在有界 $C^1$ 区域的 $W^{1,p}$ 中一致有界，用固定的有界支撑延拓得到 $f_n=Eu_n$，则 $\|f_n\|_{W^{1,p}}\le M$，所有支撑都在一个固定紧集内。令 $f_{n,\varepsilon}=\rho_\varepsilon*f_n$。对固定 $\varepsilon>0$，由 H\"older，
\[
 \|f_{n,\varepsilon}\|_\infty\le\|\rho_\varepsilon\|_{p'}M,
 \quad
 \|\nabla f_{n,\varepsilon}\|_\infty\le\|\nabla\rho_\varepsilon\|_{p'}M.
\]
这些函数支撑统一有界，且一致有界、一致 Lipschitz，因此在一致范数下可抽出收敛子列。

为看清最后一步，可用有限网格证明：给定误差 $\eta$，用有限个半径足够小的球覆盖固定支撑集，Lipschitz 界保证每个球内振荡小于 $\eta$；中心点上的函数值位于固定有界区间，将这些值以步长 $\eta$ 量化，只有有限种组合。每种组合给出一个一致范数的有限精度近似。因此该函数族完全有界，再用连续函数空间完备性即可抽取一致收敛子列。这就是此处所需的 Arzel\`a--Ascoli 论证。

\Needspace{11\baselineskip}
\subsection{Rellich 的第二步：对角子列与原序列的强收敛}

\begin{theorem}[Rellich--Kondrachov 紧嵌入]\label{thm:sobolev-rellich}
设 $\Omega$ 有界 $C^1$，$1\le p<\infty$。则
\begin{enumerate}
 \item $W^{1,p}(\Omega)\Subset L^p(\Omega)$；
 \item 若 $p<d$，则 $W^{1,p}(\Omega)\Subset L^q(\Omega)$ 对每个 $1\le q<p^*$ 成立；
 \item 若 $p=d\ge2$，则对每个有限 $q$ 紧嵌入 $L^q$；
 \item 若 $p>d$，则紧嵌入 $C^{0,\beta}(\overline\Omega)$，$0<\beta<1-d/p$。
\end{enumerate}
\end{theorem}
\begin{proof}
\textbf{先证明 $L^p$ 紧性。}取 $\varepsilon_j\downarrow0$。对 $j=1$ 抽出使 $f_{n,\varepsilon_1}$ 一致收敛的子列；再在其中抽出使 $f_{n,\varepsilon_2}$ 一致收敛的子列，依次进行。取对角子列后，对每个固定 $j$，相应光滑函数都一致收敛，从而也在固定有界支撑集上 $L^p$ 收敛。

对角子列仍记 $f_n$。给定 $\eta>0$，先选 $j$ 使 $2C_\rho M\varepsilon_j<\eta/2$，再令 $m,n$ 足够大，使光滑项的 $L^p$ 距离小于 $\eta/2$。于是
\begin{align*}
 \|f_n-f_m\|_p
 &\le\|f_n-f_{n,\varepsilon_j}\|_p
     +\|f_{n,\varepsilon_j}-f_{m,\varepsilon_j}\|_p
     +\|f_{m,\varepsilon_j}-f_m\|_p\\
 &<\eta.
\end{align*}
故原函数子列在 $L^p$ 中 Cauchy，由完备性强收敛；限制回 $\Omega$ 得第一项。

\textbf{用插值提高到次临界 $q$。}若 $p<q<p^*$，已证连续嵌入保证子列在 $L^{p^*}$ 中一致有界；选 $0<\theta<1$ 使 $1/q=\theta/p+(1-\theta)/p^*$。对差值插值，
\[
 \|u_n-u_m\|_q\le\|u_n-u_m\|_p^\theta
                    \|u_n-u_m\|_{p^*}^{1-\theta}\to0.
\]
若 $q<p$，直接用有限测度包含即可。临界 $p=d$ 时，对任意指定 $q>d$ 先选一个有限 $s>q$，用已证 $W^{1,d}\hookrightarrow L^s$，再在 $L^d$ 和 $L^s$ 间插值。

\textbf{超临界 H\"older 紧性。}Morrey 给出一致的 $C^{0,\alpha}$ 界，$\alpha=1-d/p$。在紧集 $\overline\Omega$ 上，用同样的有限网格方法得到一致收敛子列。对任意函数差 $w$，有
\[
 [w]_{C^{0,\beta}}
 \le C_{\alpha,\beta}\|w\|_\infty^{1-\beta/\alpha}
                           [w]_{C^{0,\alpha}}^{\beta/\alpha}.
\]
证明只需把 $|w(x)-w(y)|$ 同时控制为 $2\|w\|_\infty$ 与 $[w]_\alpha|x-y|^\alpha$，在两界相等的距离处分段。由于一致收敛而 $\alpha$ 半范数有界，便得 $C^{0,\beta}$ 强收敛。
\end{proof}

对任意有界开集的 $W^{1,p}_0$，同样的紧性结论可用零延拓证明，而不需要一般边界延拓。无界区域会因平移丧失紧性：取 $\phi(x-ne_1)$ 即可。临界指数处则有前面的集中反例。因此紧性需要同时防止向远处逃逸和向小尺度集中。

\subsection{零均值 Poincar\'e：用刚刚证明的紧性完成反证}

\begin{theorem}[零均值 Poincar\'e]\label{thm:sobolev-mean-poincare}
若 $\Omega$ 有界、连通且 $C^1$，$1\le p<\infty$，令 $u_\Omega=|\Omega|^{-1}\int_\Omega u$，则
\[
 \|u-u_\Omega\|_p\le C\|\nabla u\|_p.
\]
\end{theorem}
\begin{proof}
若不存在统一 $C$，则可取零均值 $v_n$，使 $\|v_n\|_p=1$、$\|\nabla v_n\|_p\to0$，方法是将反例减均值后按 $L^p$ 范数归一化。它们在 $W^{1,p}$ 中有界，由 Rellich 有子列在 $L^p$ 中强收敛于 $v$，故 $\|v\|_p=1$。另一方面测试函数恒等式取极限说明 $D_i v=0$，连通性与零梯度判据给出 $v$ 为常数。有限测度下 $L^p$ 强收敛推出积分收敛，所以 $\int v=0$，该常数只能为零，与范数为 1 矛盾。
\end{proof}
此证明明确按“零边界 Poincar\'e、独立证明 Rellich、再证明零均值 Poincar\'e”的顺序进行，没有循环依赖。

\subsection{高阶嵌入与乘积估计：以三维 \texorpdfstring{$H^2$}{H2} 为例}

高阶空间可通过逐次对导数应用一阶嵌入来理解。为避免额外引用高阶边界延拓，这里先在全空间 $\R^3$ 证明。若 $u\in H^2(\R^3)$，则 $u$ 和每个 $D_i u$ 均属于 $H^1$，因此次临界估计给出
\[
 \|u\|_6+\sum_i\|D_i u\|_6\le C\|u\|_{H^2}.
\]
于是 $u\in W^{1,6}$，再由 Morrey（$1-3/6=1/2$）得到
\begin{equation}
 H^2(\R^3)\hookrightarrow C^{0,1/2}(\R^3)\cap L^\infty(\R^3),
 \qquad\|u\|_\infty\le C\|u\|_{H^2}.
 \label{eq:sobolev-h2-embedding}
\end{equation}
同时 $\nabla u\in L^2\cap L^6$，插值得 $\|\nabla u\|_4\le C\|u\|_{H^2}$。对光滑 $u,v$，二阶乘积公式包含
\[
 D_{ij}(uv)=uD_{ij}v+vD_{ij}u+(D_i u)(D_j v)+(D_j u)(D_i v).
\]
前两项用 $L^\infty\cdot L^2$，后两项用 $L^4\cdot L^4$，连同零阶和一阶项得到
\begin{equation}
 \|uv\|_{H^2(\R^3)}\le C\|u\|_{H^2}\|v\|_{H^2}.
 \label{eq:sobolev-h2-algebra}
\end{equation}
通过全空间光滑逼近及该双线性估计可推广到一般 $H^2$ 函数。全空间的 $H^2$ 稠密性由前述截断与卷积对二阶导数重复即可证明。故三维 $H^2$ 是乘法代数，这对非线性强解估计很有用。

相比之下，三维 $H^1$ 一般不是乘法代数：取 $u=\chi(x)|x|^{-a}$，$1/4\le a<1/2$，$\chi$ 在原点附近为 1。前面的径向计算说明 $u\in H^1$，但 $u^2$ 的导数平方积分含 $\int_0^\delta r^{-4a}\dd r$，故 $u^2\notin H^1$。这明确说明何时不能直接把非线性乘积留在同一个空间。

\subsection{分数阶空间：与整数阶弱导数的关系}

对 $0<s<1$、$1\le p<\infty$，一种自然的分数阶定义是
\[
 [u]_{W^{s,p}(\Omega)}^p
 =\iint_{\Omega\times\Omega}\frac{|u(x)-u(y)|^p}{|x-y|^{d+sp}}\dd x\dd y,
 \quad\|u\|_{W^{s,p}}^p=\|u\|_p^p+[u]_{W^{s,p}}^p.
\]
它把差商在所有尺度上的平均大小作为“分数阶导数”，并非先要求经典非整数次求导。这里只给出定义和下面可直接验证的 Fourier 对应。

在 $\R^d$ 上采用酉归一化 Fourier 变换。对光滑快速衰减 $u$，令 $h=y-x$，由 Tonelli 与 Plancherel，
\begin{align*}
 [u]_{W^{s,2}(\R^d)}^2
 &=\int\frac{\|u(\cdot+h)-u\|_2^2}{|h|^{d+2s}}\dd h\\
 &=\int|\widehat u(\xi)|^2
       \left(\int\frac{|e^{i\xi\cdot h}-1|^2}{|h|^{d+2s}}\dd h\right)\dd\xi\\
 &=c_{d,s}\int|\xi|^{2s}|\widehat u(\xi)|^2\dd\xi.
\end{align*}
最后一步用旋转和缩放 $h\mapsto h/|\xi|$；常数有限，因为近原点分子为 $O(|h|^2)$，远处有界，径向积分分别是 $\int_0^1r^{1-2s}\dd r$ 和 $\int_1^\infty r^{-1-2s}\dd r$，恰好要求 $0<s<1$。

更一般定义 $H^s(\R^d)$ 的 Fourier 范数为 $\int(1+|\xi|^2)^s|\widehat u|^2$。在非负整数阶，它与弱导数的 $H^k$ 范数等价，因为 $\widehat{D^\alpha u}=(i\xi)^\alpha\widehat u$。这里使用 Fourier--Plancherel 理论作为另一路基础工具。分数阶空间解释了为何边界迹自然出现半阶正则性；一般区域上的分数阶迹和延拓不在这里证明。

\subsection{嵌入结论总表与时间紧性的边界}

\begin{center}\small
\begin{tabular}{p{0.18\linewidth}p{0.34\linewidth}p{0.35\linewidth}}
\toprule
条件 & 连续控制 & 紧性结论\\
\midrule
$1\le p<d$ & $W^{1,p}\hookrightarrow L^{p^*}$ & $L^q$ 紧，$1\le q<p^*$\\
$p=d\ge2$ & 每个有限 $L^q$ & 每个有限 $L^q$ 紧\\
$d<p<\infty$ & $C^{0,\alpha}$，$\alpha=1-d/p$ & $C^{0,\beta}$ 紧，$0<\beta<\alpha$\\
$d=1,p=1$ & 绝对连续且连续 & $L^q$ 紧，任意有限 $q$\\
\bottomrule
\end{tabular}
\end{center}
一般边界版本在本节均假设有界 $C^1$；零边界版本可采用零延拓。一维 $p=1$ 的有限 $L^q$ 紧性来自 $L^1$ 紧性、由绝对连续公式得到的统一 $L^\infty$ 界，以及插值；不宣称在 $C^0$ 中紧，因为跳变层可收缩。

空间紧性本身不能排除时间振荡，例如 $u_n(t,x)=\sin(nt)\phi(x)$。Aubin--Lions 型结论要求 $X_0\Subset X\hookrightarrow X_1$，$u_n$ 在 $L^a_tX_0$ 有界，且 $\partial_tu_n$ 在 $L^b_tX_1$ 有界（这里取 $1<a,b<\infty$），才得到 $L^a_tX$ 强收敛子列。该一般时间紧性定理在此明确作为外部定理引用；下一节会另证一个适用于常见三维 Galerkin 序列的 Hilbert 空间版本。



\section{Sobolev 工具怎样进入流体方程}\label{sec:sobolev-fluids}

\subsection{负阶空间与 Poisson 弱解：一个完整应用}

取有界区域 $\Omega$，记 $H^{-1}(\Omega)=(H^1_0(\Omega))^*$，配对记为 $\langle F,v\rangle$，定义
\[
 \|F\|_{H^{-1}}=\sup_{v\ne0}\frac{|\langle F,v\rangle|}{\|\nabla v\|_2}.
\]
这里采用 Poincar\'e 等价范数。若 $f\in L^2$，则 $\langle f,v\rangle=\int fv$ 连续，因为 $|\int fv|\le C\|f\|_2\|\nabla v\|_2$。三维中更一般有 $L^{6/5}\hookrightarrow H^{-1}$，因为 $|\int fv|\le\|f\|_{6/5}\|v\|_6\le C\|f\|_{6/5}\|\nabla v\|_2$。

考虑 $-\Delta u=F$，$u|_{\partial\Omega}=0$。对经典光滑解乘测试函数并分部积分，再推广测试空间，得到
\begin{equation}
 \int_\Omega\nabla u\cdot\nabla v\dd x=\langle F,v\rangle,
 \quad\forall v\in H^1_0(\Omega).
 \label{eq:sobolev-poisson}
\end{equation}
反过来，以此为弱解定义，不要求二阶导数在 $L^2$。

\begin{theorem}[Poisson 弱解的存在唯一性]
对每个 $F\in H^{-1}$，有唯一 $u\in H^1_0$ 满足~\eqref{eq:sobolev-poisson}，且
\[
 \|\nabla u\|_2\le\|F\|_{H^{-1}}.
\]
\end{theorem}
\begin{proof}
由 Sobolev 完备性、闭子空间性质与 Poincar\'e，$H^1_0$ 在内积 $(u,v)_V=\int\nabla u\cdot\nabla v$ 下为 Hilbert 空间。使用 Hilbert 空间的 Riesz 表示定理，连续线性泛函 $F$ 唯一表示为 $(u,\cdot)_V$，得到存在性；Riesz 定理在此作为泛函分析基础定理引用。取 $v=u$ 得 $\|\nabla u\|_2^2\le\|F\|_{H^{-1}}\|\nabla u\|_2$，即估计。两解之差取为测试函数，得到差的梯度为零，再由 Poincar\'e 得两解相同。
\end{proof}

还可以直接展开能量 $J(v)=\frac12\|\nabla v\|_2^2-\langle F,v\rangle$：弱解满足
$J(u+w)-J(u)=\frac12\|\nabla w\|_2^2$，所以是唯一极小点。这把本专题与后面变分法连接起来。更一般的有界、强制双线性型由 Lax--Milgram 定理处理。

\subsection{无散空间、压力与对流三线性型}

以下设 $\Omega$ 是有界光滑区域，$d=2$ 或 $3$，采用不可压缩、定密度、无滑移模型。令
\[
 \mathcal D_\sigma=\{\vct\varphi\in C_c^\infty(\Omega)^d:
                       \nabla\cdot\vct\varphi=0\},\qquad
 H=\overline{\mathcal D_\sigma}^{L^2},\qquad
 V=\overline{\mathcal D_\sigma}^{H^1_0}.
\]
这一定义自动给出光滑无散测试场的稠密性和 $V\Subset H\hookrightarrow V^*$。不必先借助更深的边界 Helmholtz 分解来定义这些空间。对 $\vct u,\vct v,\vct w\in V$，令
\begin{equation}
 b(\vct u,\vct v,\vct w)
 =\sum_{i,j=1}^d\int_\Omega u_i(\partial_i v_j)w_j\dd x.
 \label{eq:sobolev-trilinear}
\end{equation}
三维中用 H\"older 的 $6,2,3$ 指数及嵌入、Poincar\'e，
\begin{equation}
 |b(\vct u,\vct v,\vct w)|
 \le\|\vct u\|_6\|\nabla\vct v\|_2\|\vct w\|_3
 \le C\|\nabla\vct u\|_2\|\nabla\vct v\|_2\|\nabla\vct w\|_2.
 \label{eq:sobolev-b-bound}
\end{equation}
二维则用 $4,2,4$，同样连续。因而这个原先为光滑函数写下的非线性表达式，在 $V$ 中也确实有意义。

对光滑无散场，分部积分给出
\begin{align*}
 b(\vct u,\vct v,\vct w)+b(\vct u,\vct w,\vct v)
 &=\int_\Omega\vct u\cdot\nabla(\vct v\cdot\vct w)\dd x=0,\\
 b(\vct u,\vct v,\vct v)&=0.
\end{align*}
再由稠密性和上述连续估计推广到 $V$。这是对流项在动能估计中消失的严格空间依据。

令 $\pi=p/\rho_0$ 为运动学压力，$\nu=\mu/\rho_0>0$，并沿用全文 $\vct f$ 为单位质量力。不可压缩方程为
\[
 \partial_t\vct u+(\vct u\cdot\nabla)\vct u
 -\nu\Delta\vct u+\nabla\pi=\vct f,\qquad\nabla\cdot\vct u=0.
\]
对所有 $\vct v\in V$，其无散弱形式为
\begin{equation}
 \langle\partial_t\vct u,\vct v\rangle
 +\nu\int\nabla\vct u:\nabla\vct v
 +b(\vct u,\vct u,\vct v)=\langle\vct f,\vct v\rangle.
 \label{eq:sobolev-ns-weak}
\end{equation}
压力项因 $\int\nabla\pi\cdot\vct v=-\int\pi\nabla\cdot\vct v=0$ 而消失。恢复压力及其正则性还需 de Rham/inf--sup 等理论；不能从“测试式中没有压力”推断物理压力为零。

\subsection{能量估计：Young 如何回到流体方程}

先对光滑解或有限维 Galerkin 近似取 $\vct v=\vct u$，利用反对称性，
\[
 \frac12\frac{\mathrm d}{\mathrm dt}\|\vct u\|_2^2
 +\nu\|\nabla\vct u\|_2^2
 =\langle\vct f,\vct u\rangle
 \le\|\vct f\|_{V^*}\|\nabla\vct u\|_2.
\]
Young 取 $\varepsilon=\nu/2$ 得
\begin{equation}
 \frac12\frac{\mathrm d}{\mathrm dt}\|\vct u\|_2^2
 +\frac\nu2\|\nabla\vct u\|_2^2
 \le\frac1{2\nu}\|\vct f\|_{V^*}^2.
 \label{eq:sobolev-ns-energy}
\end{equation}
因此若 $\vct u_0\in H$、$\vct f\in L^2(0,T;V^*)$，自然的速度控制为
\begin{equation}
 \vct u\in L^\infty(0,T;H)\cap L^2(0,T;V).
 \label{eq:sobolev-energy-space}
\end{equation}
其中 $L^2(0,T;V)$ 表示强可测的 $V$ 值函数且 $\int_0^T\|\vct u(t)\|_V^2\dd t<\infty$；$L^\infty$ 表示时间上的本质有界。

对粗糙三维弱解，直接用自身作测试和写时间导数链式法则需要额外论证。经典构造先在 Galerkin 层得到等式，再以弱下半连续性得到极限的\textbf{能量不等式}。这里的计算给出先验界和构造路线，并不替代完整弱解存在性证明。

\subsection{三维的非线性项与紧性：为什么需要强收敛}

对空间中的 $\vct u\in V$，将导数移到测试场，使用 $4,4,2$ 的 H\"older 和 $L^2$--$L^6$ 插值，
\begin{align*}
 \|B(\vct u,\vct u)\|_{V^*}
 &\le C\|\vct u\|_4^2\\
 &\le C\|\vct u\|_2^{1/2}\|\vct u\|_6^{3/2}
 \le C\|\vct u\|_2^{1/2}\|\nabla\vct u\|_2^{3/2}.
\end{align*}
故三维能量界给出 $B(\vct u,\vct u)\in L^{4/3}(0,T;V^*)$；形式方程因而控制 $\partial_t\vct u$ 于同一时间空间（有限 $T$ 下 $L^2_t\subset L^{4/3}_t$）。

Galerkin 近似的导数估计需结合投影稳定性，例如使用 Stokes 特征函数基时正交投影在 $V$ 上有统一界。取得这些一致估计后，Aubin--Lions 以 $V\Subset H\hookrightarrow V^*$ 给出 $L^2(0,T;H)$ 强收敛子列。若 $\vct u_n\to\vct u$ 于时空 $L^2$，则
\[
 \|\vct u_n\otimes\vct u_n-\vct u\otimes\vct u\|_{L^1_{t,x}}
 \le(\|\vct u_n\|_{L^2_{t,x}}+\|\vct u\|_{L^2_{t,x}})
       \|\vct u_n-\vct u\|_{L^2_{t,x}}\to0.
\]
这足以使对流项对光滑测试场取极限。仅有弱收敛时，乘积一般不收敛到极限的乘积；在 $(0,2\pi)$ 上，$\sin(nx)\rightharpoonup0$ 于 $L^2$，而 $\sin^2(nx)\rightharpoonup1/2$ 于 $L^2$，已说明问题。

\subsection{二维 Ladyzhenskaya 不等式及唯一性估计}

\begin{proposition}[二维四次方估计]\label{prop:sobolev-lady}
对 $u\in H^1_0(\Omega)$、$\Omega\subset\R^2$，
\begin{equation}
 \|u\|_4^2\le C\|u\|_2\|\nabla u\|_2.
 \label{eq:sobolev-lady}
\end{equation}
\end{proposition}
\begin{proof}
先对零延拓后的光滑函数使用二维基础估计~\eqref{eq:sobolev-w11} 于 $v=u^2$：
\[
 \|u\|_4^2=\|u^2\|_2
 \le C\|\nabla(u^2)\|_1
 \le2C\|u\|_2\|\nabla u\|_2.
\]
再由 $H^1_0$ 的稠密性，对逼近序列的差使用该不等式，得到 $L^4$ 收敛并取极限。向量场通过逐分量估计得到相应版本。
\end{proof}

若两组二维能量类解 $\vct u,\vct v$ 具有相同初值和外力，令 $\vct w=\vct u-\vct v$。相应差方程的能量计算中，唯一剩下的对流项可估为
\begin{align*}
 |b(\vct w,\vct u,\vct w)|
 &\le\|\vct w\|_4^2\|\nabla\vct u\|_2\\
 &\le C\|\vct w\|_2\|\nabla\vct w\|_2\|\nabla\vct u\|_2\\
 &\le\frac\nu2\|\nabla\vct w\|_2^2
       +\frac C\nu\|\nabla\vct u\|_2^2\|\vct w\|_2^2.
\end{align*}
系数 $\|\nabla\vct u\|_2^2$ 时间可积，Gronwall 引理推出 $\vct w=0$。在严格证明中，二维估计还给出时间导数的 $L^2_tV^*$ 控制，从而可用 Hilbert 三元组链式法则正当化差方程测试；该时间链式法则在此引用。这解释了二维唯一性估计为什么能够在能量层闭合。

三维相应插值是 $\|\vct w\|_4^2\le C\|\vct w\|_2^{1/2}\|\nabla\vct w\|_2^{3/2}$。Young 后会出现 $C\nu^{-3}\|\nabla\vct u\|_2^4\|\vct w\|_2^2$；能量界只提供梯度范数平方的时间可积性，不能据此获得四次方可积性。因此这里的推导不能推出一般三维能量类弱解唯一性。

% 接在流体唯一性估计之后，补充可直接使用的 PDE 工具及证明。
\subsection{截断测试的应用：Poisson 方程的非负性}

本小节取有界 $C^1$ 区域。设 $u\in H^1_0(\Omega)$ 满足 $\int\nabla u\cdot\nabla v=\int fv$，其中 $f\in L^2$、$f\ge0$。由弱链式法则及零迹性质，$u^-=\max(-u,0)\in H^1_0$，可合法作为测试函数；并且 $\nabla u^-=-\mathbf1_{\{u<0\}}\nabla u$。因此
\[
 -\|\nabla u^-\|_2^2=\int_\Omega fu^-\dd x\ge0.
\]
左侧非正，右侧非负，故 $\nabla u^-=0$，由 Poincar\'e 得 $u^-=0$，即 $u\ge0$ a.e.。这就是不依赖经典二阶可微性的弱最大值原理的一种基本形式。

若担心“零迹经过非光滑截断仍保留”，可用前面零迹空间的证明：零延拓后负部仍在全空间 $H^1$，沿法向切片的边界值为零，故零迹刻画给出结论。更一般非齐次边界的比较原理需要相应的迹次序条件。

\subsection{差商方法怎样产生二阶弱导数：Poisson 内部正则性}

\begin{theorem}[内部 $H^2$ 正则性]\label{thm:sobolev-interior-h2}
若 $u\in H^1(\Omega)$，$f\in L^2(\Omega)$，且 $-\Delta u=f$ 于分布意义，则 $u\in H^2_{\mathrm{loc}}(\Omega)$。对 $U\Subset\Omega$，有
\[
 \|D^2u\|_{L^2(U)}\le C_{U,\Omega'}
       \big(\|f\|_{L^2(\Omega')}+\|u\|_{H^1(\Omega')}\big),
\]
其中可选固定中间子域 $U\Subset\Omega'\Subset\Omega$，常数依赖相互距离及截断。
\end{theorem}
\begin{proof}
取 $0\le\eta\le1$，$\eta\in C_c^\infty(\Omega')$，$\eta=1$ 于 $U$ 邻域。固定方向 $i$，$h$ 足够小，使所有测试与平移支撑仍在 $\Omega'$ 内；记 $z=\delta_h^iu$。用
\[
 v=-\delta_{-h}^i(\eta^2z)
\]
测试弱方程。对固定 $h$，$z\in H^1$，因此 $v\in H^1_0$ 且紧支撑；可由内部光滑逼近将分布等式推广到这个测试函数，不要求 $u$ 预先具有二阶导数。

离散分部积分 $\int a\delta_{-h}b=-\int(\delta_h a)b$ 由变量替换得到，且差商与已有的一阶弱导数交换，故
\begin{align*}
 \int\nabla u\cdot\nabla v
 &=\int\delta_h^i\nabla u\cdot\nabla(\eta^2z)\\
 &=\int\eta^2|\nabla z|^2+2\int\eta z\nabla z\cdot\nabla\eta.
\end{align*}
令 $A=\|\eta\nabla z\|_2$，$Z=\|z\|_{L^2(\operatorname{supp}\eta)}$。交叉项由 Young 控制为 $A^2/4+C_\eta Z^2$。

右侧不对 $f$ 求导，而用差商的正向估计：
\begin{align*}
 \left|\int fv\right|
 &\le\|f\|_{L^2(\Omega')}\|\delta_{-h}^i(\eta^2z)\|_2\\
 &\le\|f\|_{L^2(\Omega')}\|\partial_i(\eta^2z)\|_2\\
 &\le C_\eta\|f\|_{L^2(\Omega')}(A+Z)\\
 &\le A^2/4+C_\eta\big(\|f\|_{L^2(\Omega')}^2+Z^2\big).
\end{align*}
合并并吸收后，$A^2\le C_\eta(\|f\|_2^2+Z^2)$。另一方面差商界给出 $Z\le C\|\nabla u\|_{L^2(\Omega')}$，因为 $z$ 的积分线段仍在固定子域内。因此对每个 $j$，$\delta_h^i(D_j u)$ 在 $L^2(U)$ 中一致有界。应用差商的反向判据，得到 $D_iD_j u\in L^2(U)$ 及所述估计。
\end{proof}
这里仅证明\textbf{内部}正则性。要得全局 $H^2(\Omega)$，还须处理边界形状、边界数据和方程系数；不能因 $f\in L^2$ 就在任意粗糙区域直接断言全局二阶正则。

\subsection{无滑移边界下的 Korn 估计：对称梯度为何足够}

流体应力依赖 $D\vct u=(\nabla\vct u+\nabla\vct u^{\mathsf T})/2$，而 Sobolev 范数依赖完整梯度。对 $\vct u\in C_c^\infty(\Omega)^d$，
\[
 2\int|D\vct u|^2
 =\int|\nabla\vct u|^2+
   \sum_{i,j}\int(\partial_j u_i)(\partial_i u_j).
\]
交叉项两次分部积分给出
\[
 \sum_{i,j}\int(\partial_j u_i)(\partial_i u_j)
 =-\sum_{i,j}\int u_i\partial_i\partial_j u_j
 =\int(\nabla\cdot\vct u)^2.
\]
再由 $H^1_0$ 稠密性推广，得到恒等式
\begin{equation}
 \boxed{2\|D\vct u\|_2^2=\|\nabla\vct u\|_2^2
                +\|\nabla\cdot\vct u\|_2^2,
       \qquad\vct u\in H^1_0(\Omega)^d.}
 \label{eq:sobolev-korn-zero}
\end{equation}
因此 $\|\nabla\vct u\|_2\le\sqrt2\|D\vct u\|_2$；再配合 Poincar\'e，黏性耗散控制完整的 $H^1_0$ 范数。若还无散，则 $2\|D\vct u\|_2^2=\|\nabla\vct u\|_2^2$，与不可压缩方程的两种黏性弱形式一致。

没有零边界时，刚体转动 $\vct u(x)=Ax$、$A^{\mathsf T}=-A$ 满足 $D\vct u=0$ 而 $\nabla\vct u=A\ne0$，故不能直接使用这个估计。一般 Korn 定理须另加零阶项或去掉刚体运动；此处不证明该一般形式。

\subsection{能量类还给出哪些时空可积性}

三维中由空间插值
$\|\vct u\|_{10/3}\le\|\vct u\|_2^{2/5}\|\vct u\|_6^{3/5}$，取 $10/3$ 次幂再积分时间，
\begin{equation}
 \int_0^T\|\vct u\|_{10/3}^{10/3}\dd t
 \le C\left(\operatorname*{ess\,sup}_{0<t<T}\|\vct u(t)\|_2^{4/3}\right)
       \int_0^T\|\nabla\vct u\|_2^2\dd t.
 \label{eq:sobolev-spacetime-3d}
\end{equation}
所以三维能量类嵌入时空 $L^{10/3}((0,T)\times\Omega)$。二维则由 Ladyzhenskaya 得
\[
 \int_0^T\|\vct u\|_4^4\dd t
 \le C\left(\operatorname*{ess\,sup}_{0<t<T}\|\vct u(t)\|_2^2\right)
       \int_0^T\|\nabla\vct u\|_2^2\dd t,
\]
即二维能量类具有时空 $L^4$ 控制。这些指数由空间插值与时间积分共同决定，不能把每一时刻的 $L^6$ 控制误写成时间上的 $L^\infty_tL^6_x$ 控制。

\subsection{一个可完整证明的时间紧性版本：有限模态与高频尾部}\label{subsec:sobolev-time-compact}

下面给出前文 Aubin--Lions 用途的一个 Hilbert 空间证明。设 $V\Subset H$ 为稠密紧嵌入的 Hilbert 空间，选 $H$ 正交归一基 $e_j\in V$ 及 $\lambda_j\uparrow\infty$，使
\[
 \|v\|_H^2=\sum_j|(v,e_j)_H|^2,\qquad
 \|v\|_V^2=\sum_j\lambda_j|(v,e_j)_H|^2.
\]
这种基可由紧自伴算子的谱定理构造：在 $V$ 的内积下解 $(Kf,v)_V=(f,v)_H$，再将 $K:H\to H$ 看作紧正自伴算子。谱定理在这里作为泛函分析基础引用。流体中对应 Stokes 特征模态。

\begin{proposition}[能量界加时间导数界推出强紧性]\label{prop:sobolev-time-compact}
若 $u_n$ 在 $L^\infty(0,T;H)\cap L^2(0,T;V)$ 中一致有界，$\partial_tu_n$ 在 $L^b(0,T;V^*)$ 中一致有界，某个 $b>1$，则有子列在 $L^2(0,T;H)$ 中强收敛。
\end{proposition}
\begin{proof}
\textbf{低频模态紧。}对固定 $j$，$a_{n,j}(t)=(u_n(t),e_j)_H$ 的分布导数为 $\langle\partial_tu_n,e_j\rangle$，其 $L^b$ 范数不超过 $C\|e_j\|_V$。由一维弱导数定理，系数有绝对连续代表元，且
\[
 |a_{n,j}(t)-a_{n,j}(s)|\le C_j|t-s|^{1-1/b}.
\]
能量界还给出系数一致有界。因此对任意固定 $N$，有限个系数 $j\le N$ 可抽出一致收敛子列。沿 $N$ 取对角子列，所有有限模态同时紧。

\textbf{高频尾部一致小。}令 $P_N$ 是前 $N$ 个模态的投影，则
\[
 \int_0^T\|(I-P_N)u_n\|_H^2\dd t
 \le\lambda_{N+1}^{-1}\int_0^T\|u_n\|_V^2\dd t
 \le C\lambda_{N+1}^{-1}\longrightarrow0.
\]
给定误差，先取 $N$ 使尾部足够小，再使子列指标足够大以控制前 $N$ 个模态，便得子列在 $L^2_tH$ 中 Cauchy。Bochner $L^2$ 完备性（可按标量 $L^p$ 完备性的求和子列法证明）给出强收敛极限。
\end{proof}

三维 Galerkin 序列的能量与非线性估计提供 $b=4/3$，因此这个命题已经足够支撑前文使用的时空 $L^2$ 强收敛。使用特征基投影时，$\|P_Nv\|_V\le\|v\|_V$，这也补足了近似方程中时间导数估计所需的投影稳定性。完整的 Navier--Stokes 弱解存在性仍需核对近似初值、极限测试、初值连续性及能量不等式，本专题只把关键函数空间步骤逐一证明。

\subsection{新增练习：每题都检查一个证明环节}

\begin{enumerate}
 \item \textbf{基本引理。}为什么 $h*\rho_\varepsilon=0$ 本身还不能立刻推出 $h=0$？\\
 \textit{解答：}必须用局部 $L^1$ 近似恒等性让 $\varepsilon\to0$，并在可数穷竭子集上识别极限。
 \item \textbf{界面法向。}若把 $\Sigma$ 的法向反向，同时交换正负侧，式~\eqref{eq:weak-interface} 是否改变？\\
 \textit{解答：}$[u]$ 与 $n_i$ 同时反号，乘积不变；这是检查符号的一种方法。
 \item \textbf{差商端点。}在 $(-1,1)$ 上计算阶跃函数的 $\|\delta_hH\|_1$。\\
 \textit{解答：}充分小 $h$ 时为 1。它的差商质量集中成 Dirac 质量，不能抽出 $L^1$ 弱收敛子列来完成反向判据。
 \item \textbf{幂指数。}从 $\gamma d/(d-1)=(\gamma-1)p'$ 独立解出 $\gamma$ 和 $p^*$。\\
 \textit{解答：}依次为 $p(d-1)/(d-p)$ 与 $dp/(d-p)$；检查 $d=3,p=2$ 得 $4,6$。
 \item \textbf{Morrey 端点。}核 $|x-z|^{1-d}$ 在什么条件下属于局部 $L^{p'}$？\\
 \textit{解答：}需 $(d-1)p'<d$，等价于 $p>d$。当 $p=d$ 时径向积分出现 $\int_0^r t^{-1}\dd t$。
 \item \textbf{紧性顺序。}Rellich 证明中为何先选光滑尺度，再选序列指标？\\
 \textit{解答：}先用一致光滑误差控制整个函数族，固定尺度后才有光滑序列的收敛；颠倒顺序不能给出统一的 Cauchy 估计。
 \item \textbf{Korn。}对不可压缩无滑移场核对 $2\int|D\vct u|^2=\int|\nabla\vct u|^2$，并说明刚体转动为何不是反例。\\
 \textit{解答：}恒等式含零边界假设，非零刚体转动一般不满足无滑移条件。
 \item \textbf{时空指数。}求 $\theta$ 使 $3/10=\theta/2+(1-\theta)/6$，再核对式~\eqref{eq:sobolev-spacetime-3d}。\\
 \textit{解答：}$\theta=2/5$，取 $10/3$ 次幂后两个指数为 $4/3$ 和 $2$，恰好能由动能上界及耗散积分控制。
\end{enumerate}


\subsection{基础复习题与解答要点}

\begin{enumerate}
 \item \textbf{归一化。}不直接引用 H\"older，证明 $\int|fg|\le\|f\|_3\|g\|_{3/2}$。\\
 \textit{解答：}除以两个非零范数，应用 Young 的 $3,3/2$ 指数并积分，右侧为 $1/3+2/3=1$；零范数情形单独处理。
 \item \textbf{有限测度包含。}若 $|\Omega|=8$，$f\in L^6$，证明 $\|f\|_2\le2\|f\|_6$。\\
 \textit{解答：}$8^{1/2-1/6}=8^{1/3}=2$。说明这一结论不能推广到整个 $\R^3$。
 \item \textbf{完备性。}为何 $u_n\to u$、$Du_n\to g$ 于 $L^p$ 还必须检验一个积分恒等式？\\
 \textit{解答：}只有验证 $\int u\varphi'=-\int g\varphi$，才能确认 $g=Du$，而不是一个无关的函数极限。
 \item \textbf{折角和跳跃。}计算 $D^2|x|$ 于 $(-1,1)$，判断 $|x|$ 是否在 $H^2$ 中。\\
 \textit{解答：}$D|x|=\operatorname{sgn}x$，$D^2|x|=2\delta_0$，所以不在 $H^2$，尽管属于 $W^{1,\infty}$。
 \item \textbf{三维插值。}证明 $\|u\|_3\le\|u\|_2^{1/2}\|u\|_6^{1/2}$，并核对 $L^4$ 的指数。\\
 \textit{解答：}$1/3=(1/2)/2+(1/2)/6$；$1/4=(1/4)/2+(3/4)/6$，故 $\|u\|_4\le\|u\|_2^{1/4}\|u\|_6^{3/4}$。
 \item \textbf{边界条件。}为什么一般 $H^1$ 函数不满足 $\|u\|_2\le C\|\nabla u\|_2$？\\
 \textit{解答：}非零常数反例；需零迹、零均值或其他消除常数的条件。无滑移速度采用 $H^1_0$。
 \item \textbf{临界失紧。}三维中取 $u_n=n^{1/2}\phi(n(x-x_0))$，分别算其 $L^2,L^6$ 和梯度 $L^2$ 范数。\\
 \textit{解答：}依次为 $n^{-1}\|\phi\|_2$、$\|\phi\|_6$、$\|\nabla\phi\|_2$；这是 $H^1_0\hookrightarrow L^6$ 不紧的例子。
 \item \textbf{流体指数。}三维能量类为何控制 $\partial_t\vct u$ 于 $L^{4/3}_tV^*$ 而非自动于 $L^2_tV^*$？\\
 \textit{解答：}非线性上界含 $\|\nabla\vct u\|_2^{3/2}$，其 $4/3$ 次方才只需要梯度平方的时间积分。
\end{enumerate}

\subsection{证明范围与延伸阅读}

本专题在积分不等式与完备性之后，详细证明弱导数唯一性、一维绝对连续刻画、切片、光滑逼近、乘积与链式法则、差商判据、有界 $C^1$ 区域上的一阶延拓与迹、任意维数的次临界一阶 Sobolev 嵌入、临界有限指数嵌入、Morrey、Rellich、两种 Poincar\'e、内部 $H^2$ 正则性、零边界 Korn 及一个 Hilbert 空间时间紧性命题。Poisson 与流体估计给出相应应用。测度积分与微分基础、单位分解、Riesz/Hahn--Banach、$L^p$ 结构、紧自伴谱理论和 Fourier--Plancherel 等工具明确作为基础引用。一般 Lipschitz 边界的全部技术、分数阶迹、一般 Aubin--Lions、压力重建及完整 Navier--Stokes 弱解存在性未在这里全部证明。

可对照 John K. Hunter 的大学课程讲义
\href{https://www.math.ucdavis.edu/~hunter/pdes/pde_notes.pdf}{\textit{Notes on Partial Differential Equations}}：第 1 章复习积分空间，第 3 章系统学习弱导数、逼近、嵌入与紧性，第 4 章连接椭圆弱形式，第 6 章连接演化方程和流体模型。本文证明按上述依赖关系重新组织；课后可进一步学习 Korn 不等式、分数阶空间、Helmholtz 投影和 Stokes 算子。
